\documentclass[11pt]{article}
\usepackage{tcee}
\addbibresource{3B36.bib}
\tcetema[Tratados e instituciones de la UE]{3.B.36}{Tratados, orden jurídico e instituciones de la Unión Europea.}
\tcefecha{14/11/2026}
\begin{document}
\tceetitulo

\section*{Introducción}
\minutos{2,5}

\begin{itemize}
\item \textbf{\emph{Enganche}}:
  \begin{itemize}
  \item «Del dicho al hecho hay mucho trecho», dice el refrán. El 9 de mayo de 1950, \autor{Robert Schuman} propuso poner en común el carbón y el acero de Francia y Alemania y advirtió de que Europa no se haría de una vez ni en una obra de conjunto, sino gracias a realizaciones concretas que crearan primero una solidaridad de hecho.
  \item 76 años después, la Unión Europea tiene 27 Estados, una moneda que usan 21 de ellos, un Derecho que prima sobre las constituciones nacionales y 9 países candidatos. En junio de 2026 abrió por fin las negociaciones por capítulos con Ucrania y Moldavia\nocite{Consejo2026UcraniaMoldavia}, y en octubre la Comisión propuso cómo preparar a la Unión para crecer sin cambiar los Tratados\nocite{CE2026Ampliada}.
  \item La pregunta que recorre la exposición es: \emph{¿cómo ha convertido la Unión los dichos en hechos, y le sirven sus reglas para seguir haciéndolo con más de 30 miembros?}
  \end{itemize}

\item \textbf{\emph{Relevancia}}:
  \begin{itemize}
  \item \emph{Es la integración más profunda del mundo}: una unión aduanera, un mercado único, una moneda común y un Derecho propio que los jueces nacionales aplican como si fuera el suyo.
  \item \emph{Decide mucho de lo que hace un Estado miembro}: la política comercial, la monetaria del euro y la de competencia son suyas en exclusiva, y comparte casi todas las demás.
  \item \emph{Para España}: es miembro desde 1986 \vertema{4.A.27}, el cuarto Estado por población (10,9~\% de la Unión), con 61 de los 720 escaños del Parlamento Europeo, y en torno al 9~\% del poder de voto en el Consejo.
  \end{itemize}

\item \textbf{\emph{Contextualización}}:
  \begin{itemize}
  \item \emph{Qué es}: la \concepto{Unión Europea} (UE) es una organización de integración de 27 Estados, con personalidad jurídica desde el Tratado de Lisboa (2009), que ejerce las competencias que esos Estados le atribuyen. Es una entidad \emph{sui generis}: más que una organización internacional, porque sus normas obligan directamente a los ciudadanos y se adoptan en parte por mayoría, y menos que un Estado federal, porque solo tiene las competencias que le dan los Tratados.
  \item \emph{3 rasgos la definen como «comunidad de Derecho»}:
    \begin{itemize}
    \item un \emph{ordenamiento jurídico propio} (el \concepto{acervo}), con primacía y efecto directo;
    \item una \emph{estructura institucional} propia, con 7 instituciones independientes de los Estados;
    \item la \emph{atribución de competencias}: los Estados ceden el ejercicio de parte de su soberanía en los ámbitos que fijan los Tratados.
    \end{itemize}
  \item \emph{El objeto de esta exposición}: los Tratados, el Derecho que nace de ellos y las instituciones que lo crean y lo aplican. Las finanzas y el presupuesto \vertema{3.B.37}, el mercado interior \vertema{3.B.39}, la unión monetaria \vertema{3.B.43} y las demás políticas tienen sus propios temas.
  \end{itemize}

\item \textbf{\emph{Problemática}}:
  \begin{itemize}
  \item ¿Cómo se ha pasado de una comunidad sectorial de 6 países a una Unión de 27?
  \item ¿Qué Derecho tiene la Unión, de dónde sale su fuerza y qué competencias ejerce?
  \item ¿Qué nivel de gobierno debe ocuparse de cada cosa?
  \item ¿Quién decide en la Unión, cómo y con cuánto poder cada Estado?
  \end{itemize}

\item \textbf{\emph{Estructura}}:
  \begin{itemize}
  \item La exposición va de la historia al Derecho y del Derecho a quien lo hace:
    \begin{itemize}
    \item \concepto{I. Los tratados} (\ref{sec:tratados}): de París a Lisboa, lo que aportó cada tratado, el Brexit y el debate actual sobre la ampliación y la reforma.
    \item \concepto{II. El ordenamiento jurídico} (\ref{sec:ordenamiento}): primacía y efecto directo, fuentes, competencias y principios, con el modelo de la asignación de competencias entre niveles de gobierno, y los procedimientos legislativos.
    \item \concepto{III. Las instituciones} (\ref{sec:instituciones}): composición, funciones y forma de decidir de cada una, con el modelo del poder de voto en el Consejo.
    \end{itemize}
  \end{itemize}
\end{itemize}

\begin{esquema}[unbreakable]
\textbf{3.B.36. Tratados, orden jurídico e instituciones de la Unión Europea}\\[2pt]
Introducción\esqtiempo{2,5}\\
I. Los tratados: de París a Lisboa\esqtiempo{8}\\
\hspace*{1.5em}1. Antecedentes (1945-1951)\\
\hspace*{1.5em}2. Origen y expansión (1951-1973)\\
\hspace*{1.5em}3. Crisis (1973-1985)\\
\hspace*{1.5em}4. Relanzamiento (1985-1990)\\
\hspace*{1.5em}5. De la Comunidad a la Unión Europea (1990-2007)\\
\hspace*{1.5em}6. Base actual: Lisboa, Brexit y ampliación\\
II. El ordenamiento jurídico de la Unión Europea\esqtiempo{9}\\
\hspace*{1.5em}1. Bases del ordenamiento jurídico\\
\hspace*{1.5em}2. Competencias y principios\\
\hspace*{1.5em}3. Aprobación de la legislación\\
III. Las instituciones de la Unión Europea\esqtiempo{9,5}\\
\hspace*{1.5em}1. La Comisión Europea\\
\hspace*{1.5em}2. El Parlamento Europeo\\
\hspace*{1.5em}3. El Consejo de la Unión Europea\\
\hspace*{1.5em}4. El Consejo Europeo\\
\hspace*{1.5em}5. El Tribunal de Justicia de la Unión Europea\\
\hspace*{1.5em}6. El Tribunal de Cuentas\\
\hspace*{1.5em}7. El Banco Central Europeo\\
\hspace*{1.5em}8. Otros órganos y organismos\\
Conclusión\esqtiempo{1}
\end{esquema}

\begin{notaopositor}
Reparto del cante: introducción, \tiempo{2,5}; I, \tiempo{8} (I.1-I.4, \tiempo{3,5}; I.5, \tiempo{2}; I.6, \tiempo{2,5}); II, \tiempo{9} (II.1, \tiempo{3,5}; II.2, \tiempo{3,5}, de ellos \tiempo{2} para el modelo de la asignación de competencias; II.3, \tiempo{2}); III, \tiempo{9,5} (III.1-III.2, \tiempo{3}; III.3, \tiempo{3}, de ellos \tiempo{1,5} para el poder de voto; III.4-III.8, \tiempo{3,5}), y conclusión, \tiempo{1}. La historia se canta deprisa, tratado a tratado y con lo que aportó cada uno; el peso va al Derecho y a las instituciones. Los gráficos y figuras imprescindibles son la estructura de los tratados, el de la asignación de competencias, el procedimiento legislativo ordinario y el del poder de voto. Lo que va en cajas grises, en notas y en los anexos es para las preguntas del tribunal.
\end{notaopositor}

\begin{anotaciones}
\textbf{Notación.} En el modelo de la asignación de competencias, $G_i$ es la cantidad de un bien público que se ofrece en el Estado $i$, $a_i$ mide cuánto lo valora, $c$ es su coste marginal, $e$ la externalidad que reciben los demás Estados por cada unidad y $\Delta=a_2-a_1$ la diferencia de preferencias. En el del poder de voto, $N$ es el conjunto de los $n=27$ Estados, $T\subseteq N$ una coalición, $v(T)\in\{0,1\}$ vale 1 si la coalición gana, $\phi_i(v)$ es el índice de \autor{Shapley}-\autor{Shubik} y $\beta_i$ el de \autor{Banzhaf}. TUE y TFUE son el Tratado de la Unión Europea y el Tratado de Funcionamiento de la Unión Europea.
\end{anotaciones}

% ==========================================================================
\section{Los tratados: de París a Lisboa}\label{sec:tratados}
\minutos{8}

\begin{itemize}
\item Empezamos por la historia de los Tratados, porque cada uno respondió a un problema de su tiempo y dejó una pieza del edificio actual.
\item \emph{Conclusión anticipada}: la integración avanzó a saltos, casi siempre tras una crisis y por la vía económica, con 2 movimientos que se alternan: \concepto{profundizar} (más competencias y más decisiones por mayoría) y \concepto{ampliar} (más Estados). Lisboa ordenó el resultado; la ampliación hacia el Este vuelve a exigir reformas.
\end{itemize}

% --------------------------------------------------------------------------
\subsection{Antecedentes (1945-1951)}\label{sec:antecedentes}

\begin{itemize}
\item \emph{El contexto}: una Europa destruida por la Segunda Guerra Mundial, 2 superpotencias que se repartían el continente y unos Estados Unidos que querían una Europa occidental fuerte frente a la Unión Soviética, pero sin tutelarla para siempre.
  \begin{itemize}
  \item En lo económico, Estados Unidos impulsó un orden multilateral: el Fondo Monetario Internacional, el Banco Mundial y el GATT \vertema{3.B.20}.
  \item El \concepto{Plan Marshall} (1948)\footnote{Su nombre oficial era \emph{European Recovery Program} (Programa de Recuperación Europea). Repartió unos 13.000 millones de dólares entre 1948 y 1952 en 16 países de Europa occidental; los del Este, bajo control soviético, renunciaron a él.} condicionó la ayuda a que los europeos cooperaran entre sí y liberalizaran su comercio y sus pagos: para gestionarla nacieron la \concepto{Organización Europea de Cooperación Económica} (OECE, 1948)\footnote{La OECE se transformó en 1961 en la Organización para la Cooperación y el Desarrollo Económicos (OCDE) \vertema{3.B.19}.} y la \concepto{Unión Europea de Pagos} (1950).
  \end{itemize}
\item \emph{El Congreso de La Haya} (1948) reunió a los movimientos europeístas. Había 3 ideas de Europa:
  \begin{itemize}
  \item una federación política desde el principio (los federalistas);
  \item una cooperación entre gobiernos, sin ceder soberanía (la posición británica);
  \item una integración \emph{funcionalista}, por sectores concretos, que fuera creando intereses comunes de hecho (la de \autor{Jean Monnet}).
  \end{itemize}
\item \emph{Su resultado fue el Consejo de Europa} (Tratado de Londres, 1949)\footnote{El Consejo de Europa es una organización internacional distinta de la UE, con sede en Estrasburgo y 46 Estados; su obra principal es el Convenio Europeo de Derechos Humanos (1950), que aplica el Tribunal Europeo de Derechos Humanos. No hay que confundirlo con el Consejo Europeo ni con el Consejo de la UE, que son instituciones de la Unión (\ref{sec:otros}).}, puramente intergubernamental. Pronto se vio que, para avanzar, el método tenía que ser el funcionalista: objetivos concretos, sobre todo económicos.
\end{itemize}

\begin{anotaciones}
\textbf{¿Por qué se integran los Estados? Las teorías de la integración.} La ciencia política ofrece 3 explicaciones, útiles para las preguntas del tribunal:
\begin{itemize}
\item \emph{Neofuncionalismo} (\autor{Ernst Haas}, 1958)\nocite{Haas1958Uniting}: la integración de un sector crea presiones para integrar los vecinos (\concepto{efecto derrame}, \emph{spillover}): el carbón y el acero llevaron al mercado común, el mercado común a la moneda y la moneda a la unión bancaria. Las instituciones supranacionales y los grupos de interés empujan el proceso. Explica bien los años 50 y el relanzamiento de los 80; peor, los frenazos.
\item \emph{Intergubernamentalismo liberal} (\autor{Andrew Moravcsik}, 1998)\nocite{Moravcsik1998Choice}: los gobiernos forman sus preferencias según los intereses económicos de sus países, negocian entre ellos (pesan más los grandes) y delegan en instituciones comunes para que sus compromisos sean creíbles. Explica los grandes tratados como pactos entre Francia, Alemania y el Reino Unido.
\item \emph{Postfuncionalismo} (\autor{Liesbet Hooghe} y \autor{Gary Marks}, 2009)\nocite{HoogheMarks2009Postfunctionalist}: desde Maastricht, la integración afecta a la identidad y se politiza; el «consenso permisivo» de los ciudadanos se convierte en un «disenso restrictivo» (referendos de 2005, Brexit).
\end{itemize}
\end{anotaciones}

% --------------------------------------------------------------------------
\subsection{Origen y expansión (1951-1973)}\label{sec:origen}

\subsubsection{El Tratado de París (CECA, 1951)}\label{sec:paris}

\begin{itemize}
\item \emph{La Declaración Schuman} (9 de mayo de 1950, hoy Día de Europa): \autor{Robert Schuman}, ministro francés de Asuntos Exteriores,\footnote{No hay que confundirlo con \autor{Robert Schumann}, el compositor alemán del siglo XIX.} propuso, con la idea de \autor{Jean Monnet}, poner la producción de carbón y acero de Francia y Alemania bajo una Alta Autoridad común, abierta a otros países\nocite{MangasLinan2016Instituciones}.\footnote{Francia temía la recuperación de la industria pesada alemana, base de la industria militar. Ponerla en común hacía «materialmente imposible» una nueva guerra entre los dos países.}
\item \emph{El Tratado de París} (18 de abril de 1951; en vigor el 23 de julio de 1952) creó la \concepto{Comunidad Europea del Carbón y del Acero} (CECA) con «los Seis»: Francia, la República Federal de Alemania, Italia, Bélgica, los Países Bajos y Luxemburgo. El Reino Unido no quiso entrar. Se firmó por 50 años y expiró el 23 de julio de 2002.
\item \emph{Lo que aportó}:
  \begin{itemize}
  \item un \emph{objetivo político} (la paz) con \emph{medios económicos}: un mercado común del carbón y del acero, sin aranceles ni cuotas, y con límites a las ayudas públicas;
  \item el \emph{primer marco supranacional}: una \emph{Alta Autoridad} independiente de los gobiernos (antecedente de la Comisión), un \emph{Consejo de Ministros} (lo pidieron los países del Benelux, para equilibrar el poder de la Alta Autoridad), una \emph{Asamblea Común} (antecedente del Parlamento), un \emph{Tribunal de Justicia} y un Comité Consultivo.
  \end{itemize}
\item \emph{Un fracaso que marcó el camino}: la Comunidad Europea de Defensa, firmada en 1952, murió en 1954 al rechazarla la Asamblea Nacional francesa. La defensa quedó en la OTAN (1949) y la integración siguió por la economía.
\end{itemize}

\subsubsection{Los Tratados de Roma (CEE y Euratom, 1957)}\label{sec:roma}

\begin{itemize}
\item \emph{El camino}: la Conferencia de Mesina (1955) encargó a \autor{Paul-Henri Spaak} un informe sobre una unión económica general y otra en la energía nuclear; se aprobó en Venecia en 1956. El Reino Unido volvió a quedarse fuera y prefirió una zona de libre comercio.\footnote{En 1960 fundó con Suecia, Noruega, Dinamarca, Austria, Suiza y Portugal la Asociación Europea de Libre Comercio (AELC; EFTA en inglés); Islandia entró en 1970 y Liechtenstein en 1991 \vertema{3.B.35}.}
\item \emph{Los Tratados de Roma} (25 de marzo de 1957; en vigor el 1 de enero de 1958) crearon la \concepto{Comunidad Económica Europea} (CEE) y la \concepto{Comunidad Europea de la Energía Atómica} (Euratom).
\item \emph{Lo que aportó el Tratado CEE}, una integración ya \emph{general} (de toda la economía) y no sectorial:
  \begin{itemize}
  \item \emph{objetivos por etapas}: una \concepto{unión aduanera} en 12 años (se completó el 1 de julio de 1968, antes de plazo), un \concepto{mercado común} con libre circulación de mercancías, personas, servicios y capitales, y la aproximación de las políticas económicas \vertema{3.B.10};
  \item \emph{4 políticas comunes}: agrícola, comercial, de transportes y de competencia. La comercial y la agrícola cobraron fuerza enseguida; la de transportes y la de competencia, apenas hasta los años 80;\footnote{El Tratado añadió una \emph{cláusula de flexibilidad} (art. 235 CEE, hoy art. 352 TFUE): la Comunidad podía actuar, por unanimidad, para lograr un objetivo del Tratado aunque este no le diera los poderes necesarios. Con ella entraron, desde 1972, la política regional, la de medio ambiente o la industrial, antes de que los Tratados las recogieran.}
  \item el \emph{Fondo Social Europeo} y el \emph{Banco Europeo de Inversiones} \vertema{3.B.30};
  \item la \emph{mayoría cualificada} en el Consejo para muchas decisiones, aunque en la práctica se decidió casi todo por unanimidad.
  \end{itemize}
\item \emph{El Tratado de fusión} (Bruselas, 8 de abril de 1965; en vigor el 1 de julio de 1967) unió los ejecutivos de las 3 Comunidades en una sola Comisión y un solo Consejo (la Asamblea y el Tribunal ya eran comunes desde 1958): desde entonces se habla de las \emph{Comunidades Europeas}.
\end{itemize}

\subsubsection{La expansión de los años 60 y la primera ampliación}\label{sec:expansion}

\begin{itemize}
\item \emph{Un círculo virtuoso}: el crecimiento facilitó la unión aduanera y la unión aduanera reforzó el crecimiento. La \emph{Política Agrícola Común} se diseñó en la Conferencia de Stresa (1958) \vertema{3.B.38}.
\item \emph{La crisis de la silla vacía} (1965): Francia dejó de acudir al Consejo durante 6 meses para no aceptar el paso a la mayoría cualificada. Se cerró con el \concepto{Compromiso de Luxemburgo} (1966): cuando estuvieran en juego «intereses muy importantes» de un Estado, el Consejo buscaría el acuerdo de todos. En la práctica, hasta mediados de los 80 se decidió casi siempre por unanimidad.
\item \emph{La primera ampliación} (1 de enero de 1973): el Reino Unido, Irlanda y Dinamarca. Noruega negoció su entrada, pero la rechazó en referéndum (53~\% de noes en 1972 y de nuevo en 1994). España pidió en 1962 una asociación, pero, sin democracia, solo obtuvo un Acuerdo Preferencial (1970) \vertema{4.A.27}.
\end{itemize}

% --------------------------------------------------------------------------
\subsection{Crisis (1973-1985)}\label{sec:crisis}

\begin{itemize}
\item \emph{El primer intento de unión monetaria}: el \emph{Informe Werner} (1970) propuso una unión económica y monetaria en 10 años. El fin de la convertibilidad del dólar (agosto de 1971) y la crisis del petróleo (1973) lo enterraron; quedó la \concepto{serpiente monetaria} (1972), con márgenes de $\pm2{,}25~\%$, que se deshizo en pocos años, y después el \emph{Sistema Monetario Europeo} (1979) \vertema{3.B.43}.
\item \emph{«Euroesclerosis»}: con las crisis del petróleo, los Estados levantaron barreras no arancelarias que fragmentaron el mercado, y la Comunidad perdió terreno frente a Estados Unidos y Japón en productividad y empleo.
\item \emph{Pero hubo avances institucionales}:
  \begin{itemize}
  \item los \emph{recursos propios} (Decisión de 1970 y Tratado de Luxemburgo): los derechos de aduana, los agrícolas y, desde 1979, un porcentaje del IVA sustituyeron a las contribuciones de los Estados \vertema{3.B.37};
  \item el \emph{Consejo Europeo} (cumbre de París de 1974; primera reunión en Dublín, 1975), que reúne a los jefes de Estado o de Gobierno;
  \item el \emph{Fondo Europeo de Desarrollo Regional} (FEDER, 1975) \vertema{3.B.40};
  \item el \emph{Tribunal de Cuentas} (Tratado de Bruselas de 1975; empezó a funcionar en 1977);
  \item las primeras \emph{elecciones directas al Parlamento Europeo} (1979).
  \end{itemize}
\item \emph{La ampliación al Sur}: Grecia (1981), y España y Portugal (1 de enero de 1986, con los tratados de adhesión firmados el 12 de junio de 1985), que buscaban asentar sus democracias y modernizar sus economías \vertema{4.A.27}.
\end{itemize}

% --------------------------------------------------------------------------
\subsection{Relanzamiento (1985-1990)}\label{sec:relanzamiento}

\begin{itemize}
\item \emph{El Libro Blanco} (1985): con \autor{Jacques Delors} al frente de la Comisión, se listaron unas 280 medidas para eliminar las barreras físicas, técnicas y fiscales y completar el \concepto{mercado interior} el 31 de diciembre de 1992 \vertema{3.B.39}.
\item \emph{El Acta Única Europea} (firmada en febrero de 1986; en vigor el 1 de julio de 1987), la primera gran reforma de los Tratados:
  \begin{itemize}
  \item fijó el objetivo del mercado interior y, para lograrlo, extendió la \emph{mayoría cualificada} a la armonización de legislaciones;
  \item añadió competencias: medio ambiente, investigación, cohesión económica y social, y política social;
  \item reforzó al Parlamento con el \emph{procedimiento de cooperación} y el dictamen conforme para las adhesiones;
  \item dio base jurídica al Consejo Europeo y a la cooperación en política exterior.
  \end{itemize}
\item \emph{Lo siguiente fue la moneda}: el \emph{Informe Delors} (1989) diseñó la unión económica y monetaria en 3 etapas, y la caída del Muro y la reunificación alemana (1990) aceleraron la unión política. Se convocaron 2 conferencias intergubernamentales, una para cada unión, que acabaron en Maastricht.
\end{itemize}

% --------------------------------------------------------------------------
\subsection{De la Comunidad a la Unión Europea (1990-2007)}\label{sec:reforma}

\subsubsection{El Tratado de Maastricht (1992)}\label{sec:maastricht}

\begin{itemize}
\item \emph{El Tratado de la Unión Europea} (firmado el 7 de febrero de 1992; en vigor el 1 de noviembre de 1993) creó la \concepto{Unión Europea}, aunque sin personalidad jurídica, sobre \concepto{3 pilares} (figura~\ref{figu:tratados}):
  \begin{itemize}
  \item \emph{I. Las Comunidades Europeas}: la CEE pasó a llamarse Comunidad Europea (CE), con la CECA y Euratom. Es el pilar \emph{supranacional} (método comunitario).
  \item \emph{II. La Política Exterior y de Seguridad Común} (PESC), \emph{intergubernamental}.
  \item \emph{III. La cooperación en Justicia y Asuntos de Interior} (JAI), también intergubernamental.
  \end{itemize}
\item \emph{Lo que aportó}:
  \begin{itemize}
  \item la \emph{unión económica y monetaria}, con sus etapas, los criterios de convergencia y un Banco Central Europeo independiente \vertema{3.B.43};
  \item la \concepto{ciudadanía de la Unión}, que se suma a la nacional: libre circulación y residencia, voto en las elecciones municipales y europeas donde se reside, protección diplomática de otro Estado y petición al Parlamento y al nuevo \emph{Defensor del Pueblo Europeo};
  \item el \emph{procedimiento de codecisión} entre el Parlamento y el Consejo, más mayoría cualificada, el \emph{principio de subsidiariedad} en el Tratado, el \emph{Fondo de Cohesión} y el \emph{Comité de las Regiones};
  \item el Tribunal de Cuentas pasó a ser institución.
  \end{itemize}
\item \emph{Su ratificación fue difícil}: Dinamarca lo rechazó en un primer referéndum (1992) y lo aprobó en 1993 con excepciones; en Francia, el sí ganó por poco. Fue el fin del «consenso permisivo».
\item \emph{La cuarta ampliación} (1995): Austria, Finlandia y Suecia. Antes, en 1994, entró en vigor el \concepto{Espacio Económico Europeo} (EEE), que extiende el mercado interior a Islandia, Liechtenstein y Noruega, sin la PAC ni la unión aduanera; Suiza lo rechazó en referéndum y se relaciona con la Unión por acuerdos bilaterales \vertema{3.B.35}.
\end{itemize}

\begin{infografia}[H]
\centering
\caption{De los 3 pilares de Maastricht a los 2 Tratados de Lisboa}\label{figu:tratados}
\grafico[0.97\textwidth]{estructura-tratados}
\fuente{elaboración propia a partir del TUE y del TFUE\nocite{UE2016TUE,UE2016TFUE} y de \textcite{BaldwinWyplosz2020Economics}.}
\end{infografia}

\subsubsection{Los Tratados de Ámsterdam (1997) y Niza (2001)}\label{sec:amsterdam-niza}

\begin{itemize}
\item \emph{Ámsterdam} (firmado en octubre de 1997; en vigor el 1 de mayo de 1999) acercó la Unión a los ciudadanos, pero no hizo la reforma institucional que exigía la ampliación al Este:
  \begin{itemize}
  \item incorporó el \emph{acervo de Schengen} (el acuerdo de 1985 y su convenio de 1990, que suprimen los controles en las fronteras interiores) y pasó al pilar comunitario el asilo, los visados y la inmigración;
  \item creó la \emph{Estrategia Europea de Empleo} y el cargo de Alto Representante para la PESC;
  \item extendió la codecisión y la mayoría cualificada;
  \item introdujo la \concepto{cooperación reforzada}: un grupo de Estados puede avanzar sin los demás dentro del marco de la Unión (\ref{sec:flexibilidad}).\footnote{Es la base jurídica de la «Europa a varias velocidades», que ya existía de hecho con el euro y con Schengen: no todos los Estados avanzan al mismo ritmo ni en los mismos ámbitos.}
  \end{itemize}
\item \emph{Niza} (firmado en febrero de 2001; en vigor el 1 de febrero de 2003) hizo los cambios mínimos para la gran ampliación:
  \begin{itemize}
  \item reponderó los votos en el Consejo con una \emph{triple mayoría} (votos ponderados, mayoría de Estados y, a petición, el 62~\% de la población);
  \item limitó la Comisión a un comisario por Estado (los grandes perdieron el segundo) y reforzó a su presidente;
  \item repartió de nuevo los escaños del Parlamento y facilitó la cooperación reforzada (bastaban 8 Estados).
  \end{itemize}
\item \emph{La quinta ampliación} (1 de mayo de 2004), la mayor de todas: Chipre, Eslovaquia, Eslovenia, Estonia, Hungría, Letonia, Lituania, Malta, Polonia y la República Checa (de 15 a 25). Bulgaria y Rumanía entraron el 1 de enero de 2007 (gráfico~\ref{fig:ampliaciones}).
\end{itemize}

\subsubsection{La Constitución Europea que no fue (2004)}\label{sec:constitucion}

\begin{itemize}
\item La Declaración de Laeken (2001) convocó una \emph{Convención sobre el futuro de Europa}, presidida por \autor{Valéry Giscard d'Estaing}, que redactó un \emph{Tratado por el que se establece una Constitución para Europa}, firmado en Roma en octubre de 2004. Sustituía todos los Tratados por un texto único, con personalidad jurídica para la Unión, la doble mayoría en el Consejo y la Carta de Derechos Fundamentales vinculante.
\item \emph{Fracasó en 2005}: Francia (55~\% de noes) y los Países Bajos (62~\%) lo rechazaron en referéndum, cuando ya lo habían ratificado más de la mitad de los Estados. Pesaron los símbolos de Estado (el nombre de Constitución, la bandera, el himno) y el malestar con la ampliación y la globalización.
\end{itemize}

% --------------------------------------------------------------------------
\subsection{Base actual: Lisboa, Brexit y ampliación}\label{sec:actual}

\subsubsection{El Tratado de Lisboa (2007)}\label{sec:lisboa}

\begin{itemize}
\item Tras un «periodo de reflexión», la presidencia alemana de 2007 rescató el contenido de la Constitución sin sus símbolos. El \emph{Tratado de Lisboa} se firmó el 13 de diciembre de 2007 y entró en vigor el 1 de diciembre de 2009, después de que Irlanda, el único país que lo sometió a referéndum, lo rechazara en 2008 y lo aprobara en 2009 con garantías.
\item \emph{No es un texto nuevo, sino una reforma} de los Tratados existentes, que quedan así (figura~\ref{figu:tratados}):
  \begin{itemize}
  \item el \concepto{Tratado de la Unión Europea} (TUE): valores, principios, instituciones, cooperación reforzada, acción exterior y PESC;
  \item el \concepto{Tratado de Funcionamiento de la Unión Europea} (TFUE), el antiguo Tratado CE: las competencias, las políticas y el funcionamiento de las instituciones;
  \item la \concepto{Carta de los Derechos Fundamentales de la Unión Europea}, proclamada en Niza en 2000, con el mismo valor jurídico que los Tratados (art. 6 TUE)\nocite{UE2016Carta};\footnote{La Carta reúne en 6 capítulos (dignidad, libertades, igualdad, solidaridad, ciudadanía y justicia) los derechos civiles, políticos, económicos y sociales. Obliga a las instituciones y a los Estados cuando aplican el Derecho de la Unión. Polonia obtuvo un protocolo que limita su invocación ante sus tribunales (el Reino Unido lo tenía también).}
  \item el \emph{Tratado Euratom}, que sigue separado.
  \end{itemize}
\item \emph{Lo que aportó}:
  \begin{itemize}
  \item \emph{personalidad jurídica} única para la Unión (art. 47 TUE) y el fin de los pilares; la PESC conserva reglas propias;
  \item un \emph{reparto de competencias} claro (exclusivas, compartidas y de apoyo) y un papel para los Parlamentos nacionales en el control de la subsidiariedad (\ref{sec:competencias});
  \item el \emph{procedimiento legislativo ordinario} (la antigua codecisión) como regla general, lo que pone al Parlamento a la par del Consejo;
  \item la \emph{doble mayoría} en el Consejo desde 2014 (\ref{sec:consejo});
  \item el Consejo Europeo y el BCE como instituciones, un \emph{presidente estable del Consejo Europeo} y un \emph{Alto Representante} que es también vicepresidente de la Comisión, con un Servicio Europeo de Acción Exterior;
  \item la \emph{iniciativa ciudadana europea} y, por primera vez, una \emph{cláusula de retirada} (art. 50 TUE).
  \end{itemize}
\item \emph{La séptima ampliación}: Croacia, el 1 de julio de 2013 (28 Estados).
\end{itemize}

\subsubsection{La retirada del Reino Unido (Brexit)}\label{sec:brexit}

\begin{itemize}
\item \emph{El artículo 50 TUE} permite a un Estado retirarse «de conformidad con sus normas constitucionales»: lo notifica al Consejo Europeo, la Unión negocia un acuerdo de retirada (lo celebra el Consejo por mayoría cualificada, con la aprobación del Parlamento) y, si no hay acuerdo, los Tratados dejan de aplicarse a los 2 años, salvo prórroga por unanimidad.
\item \emph{El caso británico}:
  \begin{itemize}
  \item en el referéndum del 23 de junio de 2016, el 51,9~\% votó salir;
  \item el Reino Unido notificó la retirada el 29 de marzo de 2017; tras 3 prórrogas, salió el 31 de enero de 2020 con un \emph{Acuerdo de Retirada} (derechos de los ciudadanos, liquidación financiera y el Protocolo sobre Irlanda del Norte)\nocite{UE2020Retirada};
  \item tras un periodo transitorio hasta el 31 de diciembre de 2020, la relación se rige por el \emph{Acuerdo de Comercio y Cooperación} (en vigor desde 2021), una zona de libre comercio sin aranceles ni cuotas, pero con controles aduaneros. El Marco de Windsor (2023) ajustó el régimen de Irlanda del Norte, y la cumbre de mayo de 2025 abrió una agenda de acercamiento \vertema{3.B.41} \vertema{3.B.42}.
  \end{itemize}
\item \emph{La lección}: salir es posible, pero caro, y nadie ha querido seguir el ejemplo.
\end{itemize}

\subsubsection{Ampliación y reforma de los Tratados:\cortetitulo situación en 2026}\label{sec:ampliacion}

\begin{itemize}
\item \emph{Cómo se entra} (art. 49 TUE): el Estado europeo que respete los valores del art. 2 lo pide; el Consejo decide por unanimidad, tras consultar a la Comisión y con la aprobación del Parlamento, y el tratado de adhesión lo ratifican todos los Estados. Las condiciones son los \concepto{criterios de Copenhague} (1993): instituciones democráticas y Estado de Derecho, una economía de mercado capaz de competir en la Unión y la capacidad de asumir el acervo; a ellos se suma la capacidad de la Unión para absorber nuevos miembros.
\item \emph{La situación en noviembre de 2026} (gráfico~\ref{fig:ampliaciones}):
  \begin{itemize}
  \item hay 9 candidatos (Albania, Bosnia y Herzegovina, Georgia, Macedonia del Norte, Moldavia, Montenegro, Serbia, Turquía y Ucrania) y un candidato potencial (Kosovo);
  \item \emph{Montenegro} va en cabeza: en octubre de 2026 tenía cerrados provisionalmente 21 de los 33 capítulos y aspira a cerrarlos todos a finales de 2026; le sigue \emph{Albania}\nocite{Consejo2026Conferencias};
  \item \emph{Ucrania y Moldavia}, candidatas desde 2022, abrieron en junio de 2026 el primer bloque de capítulos (los fundamentales), cuando Hungría levantó su veto tras el cambio de gobierno de abril, y en octubre los de mercado interior y competitividad;
  \item \emph{Turquía} y \emph{Georgia} tienen las negociaciones congeladas por el retroceso democrático.
  \end{itemize}
\item \emph{El debate de la reforma}: una Unión de 30 a 35 Estados no funcionaría con la unanimidad y con un comisario por Estado. Hay 3 vías:
  \begin{itemize}
  \item \emph{Reformar los Tratados} (art. 48 TUE): el Parlamento Europeo propuso el 22 de noviembre de 2023 una reforma amplia (más mayoría cualificada, también en la PESC y en fiscalidad; una Comisión más pequeña, rebautizada «Ejecutivo»; iniciativa legislativa para el Parlamento; más competencias en salud, energía o defensa)\nocite{PE2023Reforma}. El Consejo Europeo no ha convocado la convención que exigiría.
  \item \emph{Usar lo que ya permiten los Tratados}: la Comisión propuso el 6 de octubre de 2026 («Preparar una Unión más amplia») usar las \emph{cláusulas pasarela} para votar por mayoría cualificada las sanciones, la lucha contra el fraude fiscal y la fiscalidad de la energía, recurrir más a la cooperación reforzada y añadir a los futuros tratados de adhesión \emph{salvaguardias} temporales (hasta la suspensión del voto) para quien vulnere los valores, todo ello sin reformar los Tratados\nocite{CE2026Ampliada}.
  \item \emph{Integración gradual}: dejar participar a los candidatos en partes del mercado interior y de los programas antes de la adhesión.
  \end{itemize}
\end{itemize}

\begin{figure}[H]
\centering
\caption{Estados miembros de las Comunidades Europeas y de la UE y entrada en vigor de los tratados, 1952-2026}\label{fig:ampliaciones}
\grafico[0.98\textwidth]{ampliaciones}
\fuente{elaboración propia a partir de los tratados de adhesión y de retirada (EUR-Lex)\nocite{UE2020Retirada} y del Consejo de la UE\nocite{Consejo2026Conferencias}. Códigos ISO de los países (EL, Grecia).}
\end{figure}

\begin{itemize}
\item \emph{Recapitulación}: en 70 años, la Unión ha pasado de 6 a 27 Estados y del carbón y el acero a casi todo. Cada tratado añadió competencias, mayorías y poder al Parlamento. Con este edificio delante, veamos su Derecho, que es lo que hace a la Unión distinta de cualquier otra organización.
\end{itemize}

% ==========================================================================
\section{El ordenamiento jurídico\cortetitulo{} de la Unión Europea}\label{sec:ordenamiento}
\minutos{9}

\begin{itemize}
\item Pasamos del edificio a sus reglas: qué Derecho tiene la Unión, qué competencias ejerce y cómo legisla.
\item \emph{Conclusión anticipada}: el Derecho de la Unión es un ordenamiento propio que prima sobre el nacional y que los ciudadanos pueden invocar ante sus jueces; no lo dicen los Tratados, sino el Tribunal de Justicia. La Unión solo actúa donde se le atribuyen competencias y, salvo en las exclusivas, solo si lo hace mejor que los Estados: la subsidiariedad es, en el fondo, una regla económica de asignación.
\end{itemize}

% --------------------------------------------------------------------------
\subsection{Bases del ordenamiento jurídico}\label{sec:bases}

\subsubsection{Características: autonomía, primacía y efecto directo}\label{sec:caracteristicas}

\begin{itemize}
\item La Unión es una \concepto{comunidad de Derecho}: es una creación del Derecho (los Tratados) y persigue sus objetivos con el Derecho (los actos de sus instituciones)\nocite{Borchardt2017ABC}. Su ordenamiento tiene 5 rasgos:
  \begin{itemize}
  \item \emph{Autonomía}: es un ordenamiento propio, distinto del Derecho internacional y del de los Estados, porque estos le han transferido el ejercicio de competencias. Solo el Tribunal de Justicia lo interpreta en última instancia (art. 344 TFUE).
  \item \emph{Efecto directo}: crea derechos y obligaciones para los particulares, que pueden invocarlos ante sus jueces, frente al Estado (efecto \emph{vertical}) y, a veces, frente a otros particulares (efecto \emph{horizontal}). Lo estableció la sentencia \concepto{Van Gend en Loos} (1963): la Comunidad es «un nuevo ordenamiento jurídico de Derecho internacional» cuyos sujetos son también los nacionales.
  \item \emph{Primacía}: en caso de conflicto, la norma de la Unión prevalece sobre la nacional, sea anterior o posterior y de cualquier rango, incluida la constitucional. Lo estableció \concepto{Costa contra ENEL} (1964): si cada Estado pudiera dejar sin efecto el Derecho común con una ley posterior, ese Derecho no sería común. La norma nacional contraria no es nula: el juez nacional debe \emph{dejarla sin aplicar}, sin esperar a que la deroguen (\concepto{Simmenthal}, 1978).\footnote{La primacía no figura en los Tratados: la recoge la Declaración 17 aneja al Tratado de Lisboa, que remite a la jurisprudencia. Solo opera en el ámbito de las competencias de la Unión.}
  \item \emph{Unidad}: el \concepto{acervo} (\emph{acquis}) es el conjunto de normas, sentencias y compromisos de la Unión; quien entra lo asume entero.
  \item \emph{Complejidad}: las fuentes son muchas y de naturaleza distinta (\ref{sec:fuentes}).
  \end{itemize}
\item \emph{La responsabilidad del Estado}: un Estado que incumple el Derecho de la Unión (por ejemplo, porque no transpone una directiva) debe indemnizar a los particulares perjudicados si la norma les confería derechos, la violación es suficientemente caracterizada y hay relación de causalidad (\concepto{Francovich}, 1991; \emph{Brasserie du Pêcheur} y \emph{Factortame}, 1996).
\end{itemize}

\begin{table}[H]
\centering
\caption{La jurisprudencia que construyó el ordenamiento de la Unión}\label{tab:jurisprudencia}
\small
\begin{tabular}{>{\raggedright\arraybackslash}m{4.3cm}>{\centering\arraybackslash}m{1.1cm}>{\raggedright\arraybackslash}m{9.2cm}}
\hline
\textbf{Sentencia} & \textbf{Año} & \textbf{Qué estableció}\\
\hline
\emph{Van Gend en Loos} (26/62) & 1963 & Efecto directo de los Tratados: los particulares pueden invocarlos ante el juez nacional\\
\emph{Costa contra ENEL} (6/64) & 1964 & Primacía sobre las leyes nacionales, también las posteriores\\
\emph{Internationale Handelsgesellschaft} (11/70) & 1970 & Primacía también sobre las constituciones nacionales; los derechos fundamentales, como principios generales\\
\emph{AETR} (22/70) & 1971 & Competencias exteriores implícitas: donde hay normas comunes internas, la Comunidad negocia fuera\\
\emph{Van Duyn} (41/74) & 1974 & Efecto directo de las directivas claras, precisas e incondicionales\\
\emph{Simmenthal} (106/77) & 1978 & El juez nacional inaplica por sí mismo la norma nacional contraria\\
\emph{Cassis de Dijon} (120/78) & 1979 & Reconocimiento mutuo en el mercado interior \vertema{3.B.39}\\
\emph{Les Verts} (294/83) & 1986 & El Tratado es la «carta constitucional» de una comunidad de Derecho\\
\emph{Marshall} (152/84) & 1986 & Las directivas no tienen efecto directo horizontal (entre particulares)\\
\emph{Marleasing} (C-106/89) & 1990 & Interpretación conforme: el juez interpreta el Derecho nacional a la luz de la directiva\\
\emph{Francovich} (C-6/90 y C-9/90) & 1991 & Responsabilidad del Estado por incumplir el Derecho de la Unión\\
\emph{Kadi} (C-402/05 P) & 2008 & Los actos que aplican sanciones de la ONU respetan los derechos fundamentales de la Unión\\
\hline
\end{tabular}
\fuente{elaboración propia a partir de \textcite{Borchardt2017ABC} y \textcite{MangasLinan2016Instituciones}.}
\end{table}

\begin{anotaciones}
\textbf{La primacía, discutida por los tribunales constitucionales.} Algunos tribunales nacionales aceptan la primacía, pero se reservan controlar que la Unión no se exceda de sus competencias (\emph{ultra vires}) ni toque la identidad constitucional (art. 4.2 TUE).
\begin{itemize}
\item \emph{Alemania}: el 5 de mayo de 2020, el Tribunal Constitucional federal declaró \emph{ultra vires} el programa de compras de deuda pública del BCE (PSPP), que el Tribunal de Justicia había avalado en la sentencia \emph{Weiss} (2018), y dio 3 meses al BCE para justificar su proporcionalidad. El BCE explicó su análisis a través del Bundesbank y la Comisión abrió un procedimiento de infracción contra Alemania (junio de 2021), que cerró en diciembre al comprometerse el Gobierno a respetar la primacía \vertema{3.B.43}.
\item \emph{Polonia}: en octubre de 2021, su Tribunal Constitucional declaró contrarios a la Constitución varios artículos del TUE tal como los interpreta el Tribunal de Justicia. La Comisión llevó a Polonia ante el Tribunal de Justicia en 2023.
\end{itemize}
\end{anotaciones}

\paragraph{Los valores de la Unión y su defensa}

\begin{itemize}
\item \emph{Los valores} (art. 2 TUE): dignidad humana, libertad, democracia, igualdad, Estado de Derecho y derechos humanos.
\item \emph{El artículo 7 TUE} los protege en 2 fases:
  \begin{itemize}
  \item \emph{prevención}: el Consejo, por 4/5 de sus miembros y con la aprobación del Parlamento, constata un \emph{riesgo claro} de violación grave;
  \item \emph{sanción}: el Consejo Europeo, por unanimidad (sin el Estado afectado), constata una violación \emph{grave y persistente}, y el Consejo puede suspender derechos, incluido el de voto, por mayoría cualificada.
  \end{itemize}
\item \emph{Su uso}: la Comisión activó la primera fase contra Polonia (2017; la retiró en 2024) y el Parlamento contra Hungría (2018); la unanimidad de la segunda fase la hace casi inaplicable.
\item \emph{Por eso se buscaron vías presupuestarias}: el \concepto{régimen de condicionalidad} (Reglamento 2020/2092) permite suspender fondos si las vulneraciones del Estado de Derecho afectan al presupuesto. El Tribunal de Justicia lo avaló en febrero de 2022 y se aplicó a Hungría en diciembre de 2022\nocite{UE2020Condicionalidad} \vertema{3.B.37}.
\end{itemize}

\subsubsection{Fuentes y jerarquía del Derecho de la Unión}\label{sec:fuentes}

\begin{itemize}
\item Las fuentes se ordenan en una jerarquía: cada una solo es válida si respeta las de rango superior.
  \begin{itemize}
  \item \emph{1. Derecho originario}, el que crean los Estados: el TUE, el TFUE y la Carta (con el mismo valor), el Tratado Euratom, sus protocolos y anexos, los tratados que los modificaron y los de adhesión.
  \item \emph{2. Principios generales del Derecho}, que el Tribunal deduce de las tradiciones constitucionales comunes: proporcionalidad, seguridad jurídica, confianza legítima, igualdad, derechos fundamentales. Tienen rango de Derecho originario.
  \item \emph{3. Acuerdos internacionales} que celebra la Unión (art. 218 TFUE): están por debajo de los Tratados, pero por encima del Derecho derivado. Si la materia es de competencia exclusiva, los celebra solo la Unión (un acuerdo comercial puro); si es compartida, son \emph{mixtos} y los ratifican también los Estados \vertema{3.B.41}.
  \item \emph{4. Derecho derivado}: los actos de las instituciones (\ref{sec:derivado}), que desarrollan los Tratados y están subordinados a ellos.
  \end{itemize}
\item \emph{Al margen de la jerarquía}: los convenios entre Estados miembros sobre materias ligadas a la Unión pero fuera de sus competencias, como el \emph{Tratado de Estabilidad, Coordinación y Gobernanza} (2012) o el \emph{Tratado del Mecanismo Europeo de Estabilidad} \vertema{3.B.45}.
\end{itemize}

\subsubsection{El Derecho derivado}\label{sec:derivado}

\begin{itemize}
\item \emph{Los actos jurídicos} (art. 288 TFUE) son 5: 3 vinculantes y 2 no vinculantes (tabla~\ref{tab:actos}).
  \begin{itemize}
  \item \concepto{Reglamento}: de alcance general, obligatorio en todos sus elementos y directamente aplicable en todos los Estados, sin necesidad de norma nacional. Es el instrumento de la \emph{unificación} y lo más parecido a una ley. Entra en vigor en la fecha que fija o, si no fija ninguna, a los 20 días de su publicación en el \emph{Diario Oficial}.
  \item \concepto{Directiva}: obliga a los Estados destinatarios en cuanto al \emph{resultado}, y les deja elegir la forma y los medios; hay que \emph{transponerla} en el plazo que fija. Es el instrumento de la \emph{armonización}. No tiene efecto directo por sí misma, pero lo tienen sus disposiciones claras, precisas e incondicionales si el plazo ha vencido sin transposición o con una incorrecta, y solo frente al Estado (efecto vertical).
  \item \concepto{Decisión}: obligatoria en todos sus elementos; si designa destinatarios (un Estado, una empresa), solo los obliga a ellos. La usa, por ejemplo, la Comisión para prohibir una fusión o una ayuda de Estado.
  \item \concepto{Recomendación} y \concepto{dictamen}: no vinculantes. La primera sugiere una conducta; el segundo expresa una opinión.
  \end{itemize}
\end{itemize}

\begin{table}[H]
\centering
\caption{Los actos jurídicos de la Unión (art. 288 TFUE)}\label{tab:actos}
\small
\begin{tabular}{>{\raggedright\arraybackslash}m{2.75cm}>{\raggedright\arraybackslash}m{2.75cm}>{\raggedright\arraybackslash}m{2.75cm}>{\raggedright\arraybackslash}m{2.6cm}>{\raggedright\arraybackslash}m{2.9cm}}
\hline
& \textbf{Obligatoriedad} & \textbf{Alcance} & \textbf{Aplicabilidad} & \textbf{Efecto directo}\\
\hline
\textbf{Reglamento} & En todos sus elementos & General & Directa, sin transposición & Sí, vertical y horizontal\\
\textbf{Directiva} & En cuanto al resultado & Estados destinatarios & Exige transposición & Solo vertical, si es incondicional y el plazo ha vencido\\
\textbf{Decisión} & En todos sus elementos & Sus destinatarios (si los designa) & Directa & Sí, para sus destinatarios\\
\textbf{Recomendación} & No vinculante & --- & --- & No\\
\textbf{Dictamen} & No vinculante & --- & --- & No\\
\hline
\end{tabular}
\fuente{elaboración propia a partir del art. 288 TFUE\nocite{UE2016TFUE} y de \textcite{Borchardt2017ABC}.}
\end{table}

\begin{itemize}
\item \emph{Actos legislativos y no legislativos}: lo que hace legislativo a un acto no es su forma, sino el procedimiento. Un reglamento, una directiva o una decisión adoptados por un procedimiento legislativo son \concepto{actos legislativos} (art. 289 TFUE); los demás no lo son.
\item \emph{Por debajo de los actos legislativos}, la Comisión adopta 2 tipos de actos no legislativos:
  \begin{itemize}
  \item \concepto{Actos delegados} (art. 290 TFUE): el legislador delega en la Comisión la facultad de \emph{completar o modificar elementos no esenciales} de un acto legislativo (por ejemplo, actualizar un anexo técnico). El Parlamento o el Consejo pueden oponerse en un plazo (normalmente 2 meses) o revocar la delegación.
  \item \concepto{Actos de ejecución} (art. 291 TFUE): cuando hacen falta condiciones uniformes para aplicar un acto, la Comisión (o, en casos tasados, el Consejo) adopta las normas de ejecución bajo el control de los Estados, mediante comités de representantes nacionales: es la \concepto{comitología} (Reglamento 182/2011), con 2 procedimientos, el \emph{consultivo} (el comité opina) y el \emph{de examen} (un dictamen negativo por mayoría cualificada impide adoptar el acto, que puede ir a un comité de apelación)\nocite{UE2011Comitologia}.
  \end{itemize}
\item \emph{Cuánto pesa cada uno}: cada año se adoptan entre 55 y 150 actos por el procedimiento legislativo ordinario, con picos al final de cada legislatura, pero más de 1.000 actos de ejecución y unos 150-200 delegados (gráfico~\ref{fig:actos}). Buena parte de las reglas del mercado interior se escriben en ese segundo nivel.
\item \emph{El Derecho indicativo} (\emph{soft law}): comunicaciones, directrices, libros verdes y blancos, códigos de conducta o las recomendaciones específicas por país del Semestre Europeo \vertema{3.B.45}. No obligan, pero orientan conductas y sirven para interpretar, y son el instrumento del \emph{método abierto de coordinación}.
\end{itemize}

\begin{figure}[H]
\centering
\caption{Producción normativa de la UE: actos legislativos y actos delegados y de ejecución, 2010-2025}\label{fig:actos}
\grafico[0.98\textwidth]{actos-anuales}
\fuente{elaboración propia a partir de EUR-Lex, estadísticas de actos jurídicos (actos básicos y modificativos por año de adopción)\nocite{EURLex2026Estadisticas}. En (a), los picos de 2014, 2019 y 2024 coinciden con el final de cada legislatura. Antes de 2011, EUR-Lex no clasifica los actos de ejecución.}
\end{figure}

% --------------------------------------------------------------------------
\subsection{Competencias y principios}\label{sec:competencias}

\subsubsection{La atribución de competencias (arts. 2 a 6 TFUE)}\label{sec:atribucion}

\begin{itemize}
\item Lisboa ordenó las competencias en categorías (tabla~\ref{tab:competencias}):
  \begin{itemize}
  \item \concepto{Exclusivas} (art. 3 TFUE): solo la Unión legisla; los Estados, solo si ella los habilita o para aplicar sus actos.
  \item \concepto{Compartidas} (art. 4): legislan los dos, pero los Estados solo en la medida en que la Unión no lo haya hecho (\emph{preempción}). En investigación, espacio, cooperación al desarrollo y ayuda humanitaria, la acción de la Unión no impide la de los Estados (competencias \emph{paralelas}).
  \item \concepto{De apoyo, coordinación o complemento} (art. 6): la Unión apoya a los Estados sin sustituirlos y sin armonizar sus legislaciones.
  \item Y 2 categorías especiales: la \emph{coordinación de las políticas económicas y de empleo} (art. 5) y la \emph{PESC} (art. 24 TUE), con reglas propias.
  \end{itemize}
\end{itemize}

\begin{table}[H]
\centering
\caption{Las competencias de la Unión}\label{tab:competencias}
\small
\begin{tabular}{>{\raggedright\arraybackslash}m{3.4cm}>{\raggedright\arraybackslash}m{11.6cm}}
\hline
\textbf{Tipo} & \textbf{Ámbitos}\\
\hline
\textbf{Exclusivas} (art. 3) & Unión aduanera; normas de competencia para el mercado interior; política monetaria de la zona del euro; conservación de los recursos biológicos marinos (pesca); política comercial común; y celebrar acuerdos internacionales cuando lo prevea un acto legislativo o afecten a normas comunes\\
\textbf{Compartidas} (art. 4) & Mercado interior; política social (en lo que dice el TFUE); cohesión económica, social y territorial; agricultura y pesca (salvo la conservación); medio ambiente; consumidores; transportes; redes transeuropeas; energía; espacio de libertad, seguridad y justicia; seguridad en salud pública. Paralelas: investigación, desarrollo tecnológico y espacio; cooperación al desarrollo y ayuda humanitaria\\
\textbf{Coordinación} (art. 5) & Políticas económicas (orientaciones generales), de empleo y, si la Unión lo decide, sociales\\
\textbf{PESC} (art. 24 TUE) & Política exterior y de seguridad común, incluida la definición progresiva de una defensa común\\
\textbf{Apoyo} (art. 6) & Salud humana; industria; cultura; turismo; educación, formación, juventud y deporte; protección civil; cooperación administrativa\\
\hline
\end{tabular}
\fuente{elaboración propia a partir de los arts. 2 a 6 del TFUE\nocite{UE2016TFUE}.}
\end{table}

\begin{itemize}
\item \emph{La flexibilidad}: la \emph{cláusula de flexibilidad} (art. 352 TFUE) permite actuar, por unanimidad del Consejo y con la aprobación del Parlamento, si hace falta para un objetivo de los Tratados y no hay base jurídica específica; y la jurisprudencia \emph{AETR} reconoce competencias exteriores implícitas.
\end{itemize}

\subsubsection{Los principios: atribución, subsidiariedad, proporcionalidad y cooperación leal}\label{sec:principios}

\begin{itemize}
\item El art. 5 TUE separa 2 preguntas:
  \begin{itemize}
  \item \emph{¿Tiene la Unión la competencia?} La responde el \concepto{principio de atribución}: la Unión solo actúa dentro de las competencias que le atribuyen los Tratados; las demás son de los Estados.
  \item \emph{¿Debe ejercerla, y cómo?} La responden 2 principios:
    \begin{itemize}
    \item \concepto{Subsidiariedad}: fuera de sus competencias exclusivas, la Unión solo interviene si los objetivos no pueden alcanzarlos de manera suficiente los Estados (en el nivel central, regional o local) y pueden alcanzarse mejor en la Unión, «debido a la dimensión o a los efectos de la acción pretendida».
    \item \concepto{Proporcionalidad}: el contenido y la forma de la acción no exceden de lo necesario para alcanzar los objetivos (mejor una directiva que un reglamento, mejor una recomendación que una directiva, si bastan).
    \end{itemize}
  \end{itemize}
\item \emph{El control de la subsidiariedad} (Protocolo 2): los Parlamentos nacionales tienen 8 semanas para emitir dictámenes motivados. Si los dictámenes suman un tercio de sus votos (cada Parlamento tiene 2), la Comisión debe revisar la propuesta (\concepto{tarjeta amarilla}); si suman la mayoría en el procedimiento ordinario, el legislador decide si sigue (\emph{tarjeta naranja}). Ha habido 3 tarjetas amarillas: el derecho de huelga (2012), la Fiscalía Europea (2013) y el desplazamiento de trabajadores (2016)\nocite{PE2026Subsidiariedad}.
\item \emph{La cooperación leal} (art. 4.3 TUE): la Unión y los Estados se ayudan mutuamente; los Estados adoptan todas las medidas para cumplir sus obligaciones y se abstienen de poner en peligro los objetivos de la Unión.
\item ¿Qué quiere decir «alcanzarse mejor en la Unión»? La economía tiene una respuesta precisa.
\end{itemize}

\subsubsection{¿Qué nivel de gobierno?\cortetitulo{} La teoría de la asignación de competencias}\label{sec:oates}

\paragraph{a) Idea}

\begin{itemize}
\item La subsidiariedad es una versión jurídica de una vieja pregunta del federalismo fiscal: ¿qué nivel de gobierno debe ofrecer cada bien público? \vertema{4.B.14}
\item \textcite{Oates1972Fiscal} dio la respuesta de referencia, el \concepto{teorema de la descentralización}: si las preferencias difieren entre territorios y el bien no desborda sus fronteras, es mejor que lo ofrezca cada territorio a su medida. \textcite{AlesinaAngeloniEtro2005} y \textcite{Tabellini2003Principles} lo llevaron a la Unión: centralizar tiene un \emph{beneficio}, internalizar las externalidades entre Estados (y aprovechar las economías de escala), y un \emph{coste}, imponer a todos la misma política aunque sus preferencias difieran.
\end{itemize}

\paragraph{b) Modelo}

\subparagraph{i) Supuestos}

\begin{enumerate}
\item \emph{2 Estados}, $i=1,2$, y un bien público $G_i$ que se ofrece en cada uno (seguridad, investigación, protección del medio ambiente).
\item \emph{Preferencias distintas}: el beneficio marginal del bien en el Estado $i$ es $\mathit{BMg}_i(G_i)=a_i-G_i$, con $a_2>a_1$; llamamos $\Delta=a_2-a_1$ a la \emph{heterogeneidad} y $\bar a=(a_1+a_2)/2$ a la media.
\item \emph{Coste marginal constante} $c$, igual para los 2 Estados y para la Unión.
\item \emph{Externalidad}: cada unidad de $G_i$ da un beneficio $e\geq0$ al otro Estado (la contaminación que no cruza la frontera, la defensa común).
\item \emph{2 regímenes}:
  \begin{itemize}
  \item \emph{descentralización}: cada Estado elige su $G_i$ mirando solo su bienestar;
  \item \emph{centralización}: la Unión elige \emph{una misma} cantidad $G$ para los 2 Estados, la que maximiza el bienestar conjunto. La uniformidad es el supuesto clave: refleja que una política común difícilmente se adapta a cada Estado.
  \end{itemize}
\end{enumerate}

\subparagraph{ii) Desarrollo}

\begin{itemize}
\item \emph{Paso 1: el óptimo social} (la referencia). El bienestar conjunto es $W=\sum_i\big[(a_i+e-c)\,G_i-\tfrac12G_i^2\big]$, y su óptimo:
  \begin{itemize}
  \item $G_i^{s}=a_i-c+e$: cada Estado a su medida y contando el beneficio que da al otro.
  \end{itemize}
\item \emph{Paso 2: la descentralización}. Cada Estado iguala su beneficio marginal al coste, $a_i-G_i=c$:
  \begin{itemize}
  \item $G_i^{D}=a_i-c$: se adapta a las preferencias, pero se queda corto en $e$ porque ignora la externalidad;
  \item como el bienestar es cuadrático, la pérdida de cada Estado es $\tfrac12e^2$, y la total, $L^{D}=e^2$.
  \end{itemize}
\item \emph{Paso 3: la centralización}. La Unión maximiza $W$ con $G_1=G_2=G$, $\bar a+e-c-G=0$:
  \begin{itemize}
  \item $G^{C}=\bar a-c+e$: internaliza la externalidad, pero da a los 2 Estados la cantidad pensada para la media;
  \item cada Estado se aparta de su óptimo en $\Delta/2$ (uno recibe de más y otro de menos), y la pérdida total es $L^{C}=2\cdot\tfrac12\big(\tfrac{\Delta}{2}\big)^2=\tfrac{\Delta^2}{4}$ (gráfico~\ref{fig:oates}).
  \end{itemize}
\item \emph{Paso 4: la regla de asignación}. Conviene centralizar si $L^{C}<L^{D}$:
  \funcionclave[eq:asignacion]{\text{centralizar en la Unión}\iff e>\frac{\Delta}{2}}
  es decir, si la externalidad entre Estados supera la mitad de la diferencia de sus preferencias.
\end{itemize}

\begin{figure}[H]
\centering
\caption{El coste de la uniformidad: provisión descentralizada y centralizada de un bien público}\label{fig:oates}
\grafico[0.62\textwidth]{oates-uniforme}
\fuente{elaboración propia a partir de \textcite{Oates1972Fiscal}, con el ejemplo del texto ($a_1=4$, $a_2=8$, $c=2$). Sin externalidad, la provisión uniforme $G^u$ da de más al Estado 1 y de menos al Estado 2: los triángulos son la pérdida. Con externalidad $e$, el óptimo del Estado 2 sube a $G_2^s$.}
\end{figure}

\subparagraph{iii) Implicaciones}

\begin{itemize}
\item \emph{El teorema de Oates} es el caso $e=0$: entonces $L^{D}=0<L^{C}$, y descentralizar es siempre mejor si las preferencias difieren.
\item \emph{La subsidiariedad del art. 5 TUE} tiene la misma lógica: «la dimensión o los efectos» de la acción son las economías de escala y las externalidades; el coste de centralizar es la heterogeneidad de los 27.
\item \emph{Las competencias exclusivas encajan}: la política comercial y la unión aduanera (un arancel distinto por Estado rompería el mercado único), la competencia o la moneda tienen externalidades muy grandes y preferencias relativamente parecidas (gráfico~\ref{fig:asignacion}).
\item \emph{Cuantos más Estados, más heterogeneidad}: cada ampliación sube $\Delta$, lo que justifica una Unión que centralice menos ámbitos o deje más flexibilidad (cooperación reforzada).
\end{itemize}

\begin{figure}[H]
\centering
\caption{La regla de asignación de competencias entre la UE y los Estados}\label{fig:asignacion}
\grafico[0.8\textwidth]{asignacion-competencias}
\fuente{elaboración propia a partir del modelo del texto y de \textcite{AlesinaAngeloniSchuknecht2005}. La posición de cada política es orientativa. Con un ahorro fijo $F$ al centralizar, la frontera baja: centralizar si $e^2+F>\Delta^2/4$.}
\end{figure}

\paragraph{c) Extensiones}

\begin{itemize}
\item \emph{Economías de escala}: si centralizar ahorra un coste fijo $F$ (no duplicar ejércitos, agencias o redes), centralizar compensa si $e^2+F>\Delta^2/4$, una frontera más baja.
\item \emph{El tamaño de la unión}: \textcite{AlesinaSpolaore2003Size} y \textcite{AlesinaAngeloniEtro2005} hacen endógeno el número de miembros y las competencias: una unión más grande internaliza más, pero es más heterogénea; si decide por mayoría, tiende a centralizar demasiado los ámbitos que interesan al votante mediano.
\item \emph{La movilidad}: si los ciudadanos «votan con los pies» entre jurisdicciones (\autor{Tiebout}, 1956), la descentralización revela y satisface mejor las preferencias\nocite{Tiebout1956Pure}; en la Unión, la movilidad es baja.
\item \emph{La segunda generación del federalismo}: los gobiernos no son benevolentes ni están bien informados; la descentralización acerca la decisión a quien conoce el problema, pero puede crear restricciones presupuestarias blandas\nocite{Oates1999Essay} \vertema{4.B.14}.
\end{itemize}

\paragraph{d) Evidencia empírica}

\begin{itemize}
\item \textcite{AlesinaAngeloniSchuknecht2005} midieron la actividad de la Unión por ámbitos entre 1971 y 2000: creció sobre todo donde la teoría lo justifica (comercio, mercado interior, competencia, medio ambiente), pero también en ámbitos de preferencias muy distintas y pocas externalidades (la agricultura y buena parte del gasto regional), y poco en otros donde las externalidades son grandes (defensa, política exterior).
\item Las encuestas del Eurobarómetro muestran el mismo desajuste: los ciudadanos piden más Europa en defensa, migraciones, energía o clima, justo donde hay más externalidades.
\item Las 3 tarjetas amarillas de los Parlamentos nacionales, y que ninguna haya ido más allá, muestran que el control político de la subsidiariedad es débil.
\end{itemize}

\paragraph{e) Valoración}

\begin{itemize}
\item El modelo da un criterio claro y medible para la subsidiariedad y explica la lista de competencias exclusivas.
\item Pero la asignación real responde también a la historia y a la economía política: la PAC se centralizó por un pacto entre Francia y Alemania, no por sus externalidades; y la defensa sigue siendo nacional por soberanía, aunque el modelo pida centralizarla.
\item Además, el modelo supone que la Unión decide como un planificador. En la realidad decide con unas reglas e instituciones concretas, y esas reglas reparten el poder entre los Estados: es lo que veremos en el bloque~\ref{sec:instituciones}.
\end{itemize}

% --------------------------------------------------------------------------
\subsection{Aprobación de la legislación}\label{sec:legislacion}

\subsubsection{El procedimiento legislativo ordinario}\label{sec:plo}

\begin{itemize}
\item \emph{Es la regla general desde Lisboa} (arts. 289 y 294 TFUE): el Parlamento y el Consejo adoptan juntos, en pie de igualdad, un acto propuesto por la Comisión. Se aplica en más de 80 ámbitos y el Consejo vota por mayoría cualificada (figura~\ref{figu:plo}).
\item \emph{La iniciativa} es casi un monopolio de la Comisión. Pueden pedirle que proponga el Parlamento (por mayoría de sus miembros, art. 225), el Consejo (art. 241) y un millón de ciudadanos de al menos 7 Estados con una \concepto{iniciativa ciudadana europea} (art. 11.4 TUE), pero la Comisión decide.\footnote{En algunos casos la iniciativa la tienen otros:
\begin{itemize}
\item un cuarto de los Estados, en cooperación judicial penal y policial;
\item el BCE, sobre su Estatuto;
\item el Tribunal de Justicia, para crear tribunales especializados.
\end{itemize}}
\item \emph{Las lecturas}:
  \begin{itemize}
  \item \emph{1.ª lectura}: el Parlamento aprueba su posición, sin plazo; si el Consejo la acepta, el acto queda adoptado; si no, el Consejo adopta su propia posición.
  \item \emph{2.ª lectura} (3 meses, ampliables a 4, en cada institución): el Parlamento puede aprobar la posición del Consejo, rechazarla por mayoría de sus miembros (el acto decae) o enmendarla; si el Consejo no acepta todas las enmiendas, se convoca el Comité de Conciliación.
  \item \emph{Conciliación y 3.ª lectura} (6 semanas, ampliables a 8, en cada paso): una delegación del Consejo y otra igual del Parlamento buscan un texto conjunto, que luego aprueban o rechazan las 2 instituciones sin poder enmendarlo.
  \end{itemize}
\item \emph{En la práctica, los trílogos}: reuniones informales entre el Parlamento, el Consejo y la Comisión que cierran acuerdos antes de votar. Por eso casi todos los actos se adoptan en primera lectura, y la conciliación es hoy excepcional.
\end{itemize}

\begin{infografia}[H]
\centering
\caption{El procedimiento legislativo ordinario (art. 294 TFUE)}\label{figu:plo}
\grafico[0.95\textwidth]{procedimiento-legislativo}
\fuente{elaboración propia a partir del art. 294 del TFUE\nocite{UE2016TFUE} y de \textcite{Borchardt2017ABC}.}
\end{infografia}

\subsubsection{Los procedimientos legislativos especiales}\label{sec:especiales}

\begin{itemize}
\item En los ámbitos que los Tratados reservan a los Estados, el acto lo adopta el Consejo (casi siempre por unanimidad) con una intervención menor del Parlamento:
  \begin{itemize}
  \item \emph{consulta}: el Parlamento da su opinión, que no obliga, pero sin ella el acto es nulo (fiscalidad, art. 113; armonización de legislaciones por la vía del art. 115);
  \item \emph{aprobación}: el Parlamento acepta o rechaza en bloque, sin enmendar (el marco financiero plurianual, art. 312; la cláusula de flexibilidad, art. 352; las adhesiones y los acuerdos internacionales, aunque estos no sean actos legislativos).
  \end{itemize}
\item En pocos casos lo adopta el Parlamento con la aprobación del Consejo (el estatuto de los diputados o del Defensor del Pueblo).
\item \emph{El presupuesto anual} tiene su procedimiento propio (art. 314 TFUE): la Comisión presenta el proyecto, el Consejo y el Parlamento lo leen una vez cada uno y, si discrepan, un Comité de Conciliación tiene 21 días; el Parlamento tiene la última palabra en algunos supuestos \vertema{3.B.37}.
\end{itemize}

\subsubsection{La flexibilidad: cooperación reforzada y cláusulas pasarela}\label{sec:flexibilidad}

\begin{itemize}
\item \emph{La cooperación reforzada} (art. 20 TUE y arts. 326-334 TFUE): si la Unión en su conjunto no puede alcanzar un objetivo en un plazo razonable, al menos 9 Estados pueden avanzar solos dentro de sus instituciones, como último recurso, abierta a los demás y con la autorización del Consejo por mayoría cualificada (por unanimidad en la PESC). Se ha usado para la ley aplicable al divorcio (2010), la patente unitaria (2012; el Tribunal Unificado de Patentes funciona desde 2023), los regímenes económicos matrimoniales (2016) y la Fiscalía Europea (2017).
\item \emph{La cooperación estructurada permanente} (CEP o PESCO, 2017), en defensa, con 26 Estados (todos menos Malta).
\item \emph{Las cláusulas pasarela}: la \emph{general} (art. 48.7 TUE) permite al Consejo Europeo, por unanimidad, pasar un ámbito de la unanimidad a la mayoría cualificada, o de un procedimiento especial al ordinario, si ningún Parlamento nacional se opone en 6 meses; hay también pasarelas específicas (marco financiero, art. 312.2; política social, art. 153.2; medio ambiente, art. 192.2). Nunca se ha usado la general.
\item \emph{El resultado es una Europa de geometría variable}: no todos los Estados están en todo (figura~\ref{figu:circulos}).
\end{itemize}

\begin{infografia}[H]
\centering
\caption{La Europa de geometría variable en 2026}\label{figu:circulos}
\grafico[0.97\textwidth]{circulos-europa}
\fuente{elaboración propia a partir de la Unión Europea\nocite{UE2026Instituciones}, el Consejo de Europa y la AELC. Códigos ISO de los países (EL, Grecia; XK, Kosovo).}
\end{infografia}

\begin{itemize}
\item \emph{Recapitulación}: la Unión tiene un Derecho propio, que prima y se aplica directamente, dentro de las competencias que le dan los Tratados y con la subsidiariedad como regla de reparto. Ese Derecho lo hacen unas instituciones concretas, que reparten el poder entre Estados, ciudadanos e interés común: es lo que estudiamos a continuación.
\end{itemize}

% ==========================================================================
\section{Las instituciones\cortetitulo{} de la Unión Europea}\label{sec:instituciones}
\minutos{9,5}

\begin{itemize}
\item Terminamos con quién decide. El art. 13 TUE enumera 7 \concepto{instituciones}: el Parlamento Europeo, el Consejo Europeo, el Consejo, la Comisión Europea, el Tribunal de Justicia de la Unión Europea, el Banco Central Europeo y el Tribunal de Cuentas. Las completan órganos consultivos, el BEI y varios organismos y agencias.\footnote{Los tratados fundacionales tenían 4 instituciones (Comisión, Consejo, Asamblea y Tribunal de Justicia). Maastricht añadió el Tribunal de Cuentas, y Lisboa, el Consejo Europeo y el BCE.}
\item \emph{Conclusión anticipada}: no hay una correspondencia simple con la división de poderes de \autor{Montesquieu}. El poder legislativo lo comparten el Parlamento (los ciudadanos) y el Consejo (los Estados), con la iniciativa de la Comisión (el interés común); el ejecutivo se reparte entre la Comisión y el Consejo, y el judicial, entre el Tribunal de Justicia y los jueces nacionales. El resultado es un \concepto{equilibrio institucional} entre lo supranacional y lo intergubernamental (figura~\ref{figu:instituciones}).
\end{itemize}

\begin{infografia}[H]
\centering
\caption{Las instituciones de la UE: quién elige a quién y quién legisla}\label{figu:instituciones}
\grafico[0.95\textwidth]{mapa-instituciones}
\fuente{elaboración propia a partir del TUE y del TFUE\nocite{UE2016TUE,UE2016TFUE}.}
\end{infografia}

% --------------------------------------------------------------------------
\subsection{La Comisión Europea}\label{sec:comision}

\begin{itemize}
\item \emph{Qué es} (art. 17 TUE): la institución que promueve el \emph{interés general} de la Unión, con independencia de los gobiernos. Es el motor de la integración y la guardiana de los Tratados. Tiene su sede en Bruselas.\footnote{En el edificio Berlaymont. Su antecedente es la Alta Autoridad de la CECA (1952); las Comisiones de la CEE y de Euratom (1958) se fusionaron en 1967.}
\item \emph{Composición}: un \emph{Colegio} de 27 miembros, uno por Estado, con un mandato de 5 años.\footnote{Lisboa preveía reducirla a dos tercios del número de Estados desde 2014, pero el Consejo Europeo, usando la facultad del art. 17.5 TUE, mantuvo un comisario por Estado (lo había prometido a Irlanda para su segundo referéndum).}
  \begin{itemize}
  \item el \emph{presidente} lo propone el Consejo Europeo por mayoría cualificada, «teniendo en cuenta» las elecciones al Parlamento, y lo elige el Parlamento por mayoría de sus miembros;
  \item los demás comisarios los propone el Consejo de común acuerdo con el presidente electo, a partir de las propuestas de los Estados; el Parlamento examina a cada uno en audiencias y vota al Colegio entero, que nombra el Consejo Europeo por mayoría cualificada;
  \item el Colegio actual, el segundo de \autor{Ursula von der Leyen} (reelegida en julio de 2024 con 401 votos), empezó el 1 de diciembre de 2024. Tiene 6 vicepresidentes ejecutivos, entre ellos la española \autor{Teresa Ribera} (Transición Limpia, Justa y Competitiva) y la Alta Representante, \autor{Kaja Kallas}\nocite{CE2024Colegio};
  \item la administración se organiza en direcciones generales.
  \end{itemize}
\item \emph{Funciones}:
  \begin{itemize}
  \item \emph{Iniciativa legislativa}: casi siempre exclusiva; prepara sus propuestas con comunicaciones, libros verdes y blancos y evaluaciones de impacto.
  \item \emph{Ejecutiva}: ejecuta el presupuesto y gestiona los programas (directamente o con los Estados, en gestión compartida) \vertema{3.B.37}; adopta actos delegados y de ejecución, y delega tareas en agencias ejecutivas.
  \item \emph{Guardiana de los Tratados}: abre \emph{procedimientos de infracción} (carta de emplazamiento, dictamen motivado y recurso ante el Tribunal, que puede imponer multas coercitivas; arts. 258-260 TFUE) y aplica directamente la política de competencia y de ayudas de Estado \vertema{3.B.39}.
  \item \emph{Representación exterior}: salvo en la PESC, representa a la Unión y negocia los acuerdos internacionales, en particular los comerciales, con un mandato del Consejo \vertema{3.B.41}.
  \end{itemize}
\item \emph{Funcionamiento}: decide de forma colegiada, por mayoría simple si hay votación, en una reunión semanal (en Estrasburgo en las semanas de pleno) o por procedimiento escrito. Responde ante el Parlamento, que puede censurarla.
\end{itemize}

% --------------------------------------------------------------------------
\subsection{El Parlamento Europeo}\label{sec:parlamento}

\begin{itemize}
\item \emph{Qué es} (art. 14 TUE): la institución que representa directamente a los ciudadanos. Es el fundamento democrático de la Unión. Tiene 3 sedes: los plenos en Estrasburgo, las comisiones y los plenos adicionales en Bruselas y la Secretaría General en Luxemburgo.
\item \emph{Composición}:
  \begin{itemize}
  \item como máximo 750 diputados más el presidente; en la legislatura 2024-2029, \emph{720}, elegidos por sufragio universal directo cada 5 años con las leyes electorales de cada Estado (siempre proporcionales);
  \item el reparto es \concepto{decrecientemente proporcional}: entre 6 y 96 escaños por Estado, de modo que un diputado alemán representa a unos 870.000 habitantes y uno maltés a menos de 100.000. España tiene 61 (gráfico~\ref{fig:escanos})\nocite{ConsejoEuropeo2023Composicion};
  \item los diputados no se agrupan por países, sino en \emph{grupos políticos}. Ninguno tiene mayoría: la que eligió a la Comisión la forman el Partido Popular Europeo, los Socialistas y Demócratas y Renovar Europa (gráfico~\ref{fig:grupos}).\footnote{Desde 2024, los cambios de grupo han movido las cifras: en julio de 2026, el PPE tenía 184 escaños, los Socialistas y Demócratas 135, Patriotas por Europa 85, los Conservadores y Reformistas 84 y Renovar Europa 78\nocite{SIEPS2026Centre}.} La presidenta es \autor{Roberta Metsola}, reelegida en julio de 2024 para la primera mitad de la legislatura.
  \end{itemize}
\end{itemize}

\begin{figure}[H]
\centering
\caption{Escaños por Estado miembro en el Parlamento Europeo (2024-2029) y habitantes por escaño}\label{fig:escanos}
\grafico[0.8\textwidth]{escanos-estados}
\fuente{elaboración propia a partir de la Decisión (UE) 2023/2061 del Consejo Europeo\nocite{ConsejoEuropeo2023Composicion} y de las cifras de población de 2026 de la Decisión 2025/2427 del Consejo\nocite{Consejo2025Poblacion}.}
\end{figure}

\begin{figure}[H]
\centering
\caption{Composición del Parlamento Europeo por grupos políticos al iniciarse la legislatura 2024-2029}\label{fig:grupos}
\grafico[0.8\textwidth]{grupos-pe}
\fuente{elaboración propia a partir del Parlamento Europeo (sesión constitutiva del 16 de julio de 2024)\nocite{PE2024Constitutiva}. Izq.: La Izquierda; S\&D: Socialistas y Demócratas; Renew: Renovar Europa; PPE: Partido Popular Europeo; ECR: Conservadores y Reformistas Europeos; PfE: Patriotas por Europa; ESN: Europa de las Naciones Soberanas.}
\end{figure}

\begin{itemize}
\item \emph{Funciones}:
  \begin{itemize}
  \item \emph{Legislativa}: colegisla con el Consejo en el procedimiento ordinario; da su aprobación en los especiales que la exigen (adhesiones, acuerdos internacionales, marco financiero plurianual) y su opinión en la consulta. Su poder ha crecido con cada tratado, de lo consultivo a la codecisión.
  \item \emph{Presupuestaria}: aprueba con el Consejo el presupuesto anual y aprueba la gestión de la Comisión (\emph{descargo}) \vertema{3.B.37}.
  \item \emph{De control}: elige al presidente de la Comisión y aprueba el Colegio; puede aprobar una \emph{moción de censura} contra la Comisión por mayoría de dos tercios de los votos emitidos que sea a la vez mayoría de sus miembros (nunca ha prosperado; en 1999 la Comisión Santer dimitió antes de votarse); hace preguntas, crea comisiones de investigación, recibe peticiones y elige al Defensor del Pueblo Europeo.
  \end{itemize}
\end{itemize}

% --------------------------------------------------------------------------
\subsection{El Consejo de la Unión Europea}\label{sec:consejo}

\subsubsection{Composición, funciones y reglas de voto}\label{sec:consejo-reglas}

\begin{itemize}
\item \emph{Qué es} (art. 16 TUE): la institución que representa a los \emph{gobiernos} de los Estados. Lo forma un representante de cada Estado de rango ministerial, facultado para comprometer a su gobierno; tiene su sede en Bruselas (algunas reuniones, en Luxemburgo).
\item \emph{Formaciones}: el Consejo es uno, pero se reúne en 10 formaciones según el tema. Destacan la de \emph{Asuntos Generales}, que prepara el Consejo Europeo; la de \emph{Asuntos Exteriores}, que preside la Alta Representante, y la de \emph{Asuntos Económicos y Financieros} (Ecofin). Los ministros de la zona del euro se reúnen antes en el \emph{Eurogrupo}, que es informal (\ref{sec:otros}).
\item \emph{La presidencia} rota cada 6 meses entre los Estados, en \emph{tríos} con un programa común de 18 meses\nocite{Consejo2026Presidencias}:
  \begin{itemize}
  \item 2025-2026: Polonia, Dinamarca y Chipre (enero-junio de 2026);
  \item julio de 2026-diciembre de 2027: Irlanda (julio-diciembre de 2026), Lituania y Grecia;
  \item 2028-2029: Italia, Letonia y Luxemburgo. España presidió el Consejo en el segundo semestre de 2023.
  \end{itemize}
\item \emph{La preparación}: el \concepto{Comité de Representantes Permanentes} (Coreper) reúne a los embajadores de los Estados ante la Unión (Coreper II) y a sus adjuntos (Coreper I), y prepara todas las decisiones del Consejo con unos 150 grupos de trabajo; lo que llega acordado se aprueba sin debate (puntos «A»).\footnote{En la Representación Permanente de España ante la UE trabajan muchos técnicos comerciales y economistas del Estado.}
\item \emph{Funciones}: legisla con el Parlamento y aprueba el presupuesto con él; coordina las políticas económicas (Ecofin) \vertema{3.B.45}; define y aplica la PESC con las orientaciones del Consejo Europeo; y celebra los acuerdos internacionales que negocia la Comisión.
\item \emph{Reglas de voto}:
  \begin{itemize}
  \item \concepto{Mayoría cualificada} (art. 16.4 TUE), la regla general desde el 1 de noviembre de 2014, con \emph{doble mayoría}: al menos el 55~\% de los Estados (15 de 27) que representen al menos el 65~\% de la población. Si la propuesta no viene de la Comisión ni de la Alta Representante, hace falta el 72~\% de los Estados (20). La \emph{minoría de bloqueo} necesita al menos 4 Estados: así, 3 grandes que juntos superan el 35~\% de la población no bloquean solos\nocite{Consejo2026Votacion}. Las cifras de población se actualizan cada año (para 2026, el umbral es de 293,9 millones de 452,2)\nocite{Consejo2025Poblacion}.
  \item \emph{Unanimidad} (la abstención no la impide): PESC, fiscalidad, recursos propios y marco financiero, adhesiones, ciudadanía, seguridad social y protección social, cláusula de flexibilidad y parte de la justicia penal.
  \item \emph{Mayoría simple} (14 Estados): cuestiones de procedimiento.
  \end{itemize}
\end{itemize}

\subsubsection{¿Cuánto poder tiene cada Estado?\cortetitulo{} Los índices de poder de voto}\label{sec:poder}

\paragraph{a) Idea}

\begin{itemize}
\item Con la doble mayoría, ¿pesa Alemania lo que dice su población (18,5~\%)? ¿Y Malta (0,13~\%)? El \emph{peso} de un Estado no es su \emph{poder}: lo que importa es con qué frecuencia su voto decide el resultado. Un Estado con mucho peso puede no decidir nunca, y uno pequeño puede ser decisivo a menudo.
\item La teoría de juegos cooperativos mide ese poder \vertema{3.A.14}. \textcite{ShapleyShubik1954Method} aplicaron el valor de \autor{Shapley}\footnote{\autor{Lloyd Shapley} recibió el Premio Nobel de Economía en 2012, con \autor{Alvin Roth}, «por la teoría de las asignaciones estables y la práctica del diseño de mercados»\nocite{Nobel2012RothShapley}.} a las votaciones, y \textcite{Banzhaf1965Weighted} (antes, \textcite{Penrose1946Elementary}) propuso una medida parecida. Son los índices con los que se negociaron las reglas de voto de Niza y Lisboa \vertema{3.A.24} \vertema{4.B.2}.
\end{itemize}

\paragraph{b) Modelo}

\subparagraph{i) Supuestos}

\begin{enumerate}
\item \emph{Votantes}: los $n=27$ Estados, $N=\{1,\ldots,27\}$, que votan sí o no a una propuesta.
\item \emph{Regla de decisión}: una coalición $T\subseteq N$ gana, $v(T)=1$, si reúne al menos 15 Estados con al menos el 65~\% de la población, o al menos 24 Estados (porque la minoría de bloqueo exige 4); si no, $v(T)=0$. Es un \emph{juego simple}: monótono (si $T$ gana, cualquier coalición que la contenga gana).
\item \emph{Poder a priori}: no se tienen en cuenta las preferencias de los Estados; se supone que cualquier orden de llegada a la coalición (\autor{Shapley}-\autor{Shubik}) o cualquier coalición (\autor{Banzhaf}) es igual de probable. Se mide el poder que da la \emph{regla}, no la política.
\end{enumerate}

\subparagraph{ii) Desarrollo}

\begin{itemize}
\item \emph{Paso 1: el voto decisivo}. Un Estado $i$ es \emph{decisivo} (o pivote) en una coalición $T$ que no lo incluye si $T$ pierde sin él y gana con él: $v(T\cup\{i\})-v(T)=1$.
\item \emph{Paso 2: el índice de Shapley-Shubik}. Los Estados se suman a la coalición en un orden aleatorio; el poder de $i$ es la probabilidad de que sea él quien la convierta en ganadora:
  \funcionclave[eq:shapley-shubik]{\phi_i(v)=\sum_{T\subseteq N\setminus\{i\}}\frac{|T|!\,\big(n-|T|-1\big)!}{n!}\,\Big[v\big(T\cup\{i\}\big)-v(T)\Big]}
  Los índices de todos los Estados suman 1.
\item \emph{Paso 3: el índice de Banzhaf}. Todas las coaliciones son igual de probables: el poder absoluto de $i$ es la proporción de coaliciones de los demás en las que es decisivo, $\beta_i'=\frac{1}{2^{n-1}}\sum_{T\subseteq N\setminus\{i\}}\big[v(T\cup\{i\})-v(T)\big]$, y el normalizado, $\beta_i=\beta_i'/\sum_j\beta_j'$.
\item \emph{Paso 4: el cálculo para los 27}, con las cifras de población de 2026 (hay $2^{26}\approx67$ millones de coaliciones por Estado; el método, en el anexo~\ref{anexo:poder}):
  \begin{itemize}
  \item Alemania: población del 18,5~\%, \autor{Shapley}-\autor{Shubik} del 17,5~\% y \autor{Banzhaf} del 12,0~\%;
  \item España: 10,9, 9,4 y 7,8~\%; Polonia: 8,3, 6,7 y 6,3~\%;
  \item Malta: 0,13, 0,87 y 1,73~\% (gráfico~\ref{fig:poder}).
  \end{itemize}
\end{itemize}

\begin{figure}[H]
\centering
\caption{Poder de voto a priori de cada Estado en el Consejo frente a su peso en la población (2026)}\label{fig:poder}
\grafico[0.78\textwidth]{poder-voto}
\fuente{elaboración propia con la regla de la doble mayoría (arts. 16.4 TUE y 238.3 TFUE) y las cifras de población de 2026 de la Decisión 2025/2427 del Consejo\nocite{Consejo2025Poblacion}. Por encima de la bisectriz, los Estados con más poder que población.}
\end{figure}

\subparagraph{iii) Implicaciones}

\begin{itemize}
\item \emph{La doble mayoría combina 2 legitimidades}: el criterio de Estados (un Estado, un voto) favorece a los pequeños y el de población, a los grandes. El resultado da a los grandes algo menos poder que población y a los pequeños varias veces más.
\item \emph{Los 2 índices no coinciden}: \autor{Banzhaf} da más poder relativo a los pequeños que \autor{Shapley}-\autor{Shubik}, y las conclusiones sobre quién gana o pierde con una reforma dependen del índice.
\item \emph{La capacidad de decidir}: la probabilidad de que una coalición al azar sea ganadora (el poder de la colectividad para actuar, de \autor{Coleman}) es del 13,2~\% con las reglas actuales, frente a cerca del 2~\% con las de Niza en una Unión de 27\nocite{BaldwinWidgren2004Winners}: la doble mayoría hizo al Consejo mucho más ágil.
\item \emph{España}: con Niza tenía casi los votos de Alemania (27 frente a 29) con la mitad de población; con la doble mayoría perdió poder relativo frente a Alemania, aunque ganó capacidad de decidir el conjunto.
\end{itemize}

\paragraph{c) Extensiones}

\begin{itemize}
\item \emph{La ley de la raíz cuadrada} de \textcite{Penrose1946Elementary}: para que cada ciudadano tenga el mismo poder indirecto, el peso de cada Estado debe ser proporcional a la \emph{raíz cuadrada} de su población. Polonia defendió en 2007 un sistema así (el «compromiso de Jagellón», con un umbral del 61,6~\%)\nocite{SlomczynskiZyczkowski2006Penrose}.
\item \emph{El poder a posteriori}: si las preferencias de los Estados están correlacionadas (por ejemplo, frugales frente a meridionales), las coaliciones no son aleatorias, y el poder efectivo depende de quién está cerca del centro del espacio político y de quién fija la agenda (la Comisión)\nocite{FelsenthalMachover1998Measurement}.
\item \emph{La ampliación}: cada nuevo Estado diluye el poder de los demás y cambia el equilibrio entre grandes y pequeños; la entrada de un Estado grande como Ucrania desplazaría a los actuales.
\end{itemize}

\paragraph{d) Evidencia empírica}

\begin{itemize}
\item \emph{El Consejo vota poco}: aunque en torno al 80~\% de la legislación se decide por mayoría cualificada\nocite{Consejo2026Votacion}, la mayoría de los actos se adoptan sin votos en contra, porque la presidencia busca el consenso y la amenaza de la votación basta para que los reticentes cedan \nocite{HayesRenshaw2006When}(\autor{Hayes-Renshaw}, \autor{Van Aken} y \autor{Wallace}, 2006).
\item \emph{Pero el poder a priori importa en la negociación}: las negociaciones de Niza (2000) y de la Constitución (2003-2004) giraron en torno a estos índices, y España y Polonia, que eran los que más perdían, bloquearon la doble mayoría hasta que se aplazó a 2014\nocite{BaldwinWidgren2004Winners}.
\end{itemize}

\paragraph{e) Valoración}

\begin{itemize}
\item Los índices son útiles para \emph{diseñar} reglas, como tras un velo de ignorancia, y muestran que repartir votos en proporción a la población no reparte el poder en esa proporción.
\item Ignoran las preferencias, la agenda y el Parlamento como colegislador, y suponen que se vota cuando casi siempre se pacta. El poder real en la Unión está también en las instituciones que no representan a los Estados.
\end{itemize}

% --------------------------------------------------------------------------
\subsection{El Consejo Europeo}\label{sec:consejo-europeo}

\begin{itemize}
\item \emph{Qué es} (art. 15 TUE): la institución que da a la Unión el \emph{impulso político} y fija sus orientaciones y prioridades. No legisla. Nació como cumbres informales (desde 1961), se institucionalizó en 1974, entró en los Tratados con el Acta Única y es institución desde Lisboa.
\item \emph{Composición}: los jefes de Estado o de Gobierno de los 27, su presidente y la presidenta de la Comisión; participa la Alta Representante. Se reúne al menos 2 veces por semestre en Bruselas y decide por consenso, salvo que los Tratados digan otra cosa.
\item \emph{Su presidente} lo elige el propio Consejo Europeo por mayoría cualificada para 2 años y medio, renovable una vez, y no puede ejercer un mandato nacional. Desde el 1 de diciembre de 2024 es \autor{António Costa}, hasta el 31 de mayo de 2027\nocite{ConsejoEuropeo2024Presidente}.
\item \emph{Funciones}:
  \begin{itemize}
  \item \emph{orientaciones políticas} y agenda estratégica, y el desbloqueo de los grandes asuntos (el acuerdo político sobre el marco financiero plurianual, las grandes respuestas a las crisis);
  \item \emph{acción exterior}: fija los intereses y objetivos estratégicos y la PESC;
  \item \emph{gobernanza económica}: respalda las orientaciones del Semestre Europeo \vertema{3.B.45};
  \item \emph{decisiones institucionales}: la composición del Parlamento, la rotación de las presidencias, las cláusulas pasarela o la sanción del art. 7;
  \item \emph{nombramientos}: su presidente, el candidato a presidir la Comisión, el Colegio, la Alta Representante y el Comité Ejecutivo del BCE.
  \end{itemize}
\end{itemize}

% --------------------------------------------------------------------------
\subsection{El Tribunal de Justicia de la Unión Europea}\label{sec:tjue}

\begin{itemize}
\item \emph{Qué es} (art. 19 TUE): la institución que garantiza el respeto del Derecho en la interpretación y aplicación de los Tratados. Tiene su sede en Luxemburgo y nació con la CECA (1952).
\item \emph{Composición}: 2 órganos.
  \begin{itemize}
  \item El \emph{Tribunal de Justicia}: 27 jueces (uno por Estado) y 11 abogados generales, que proponen soluciones a los asuntos. Juzga en salas de 3 o 5 jueces, en Gran Sala de 15 o en Pleno.
  \item El \emph{Tribunal General} (antes, de Primera Instancia, creado en 1988): 54 jueces (2 por Estado desde 2019). Juzga en primera instancia los recursos de particulares y empresas (competencia, ayudas, marcas) y, desde octubre de 2024, las cuestiones prejudiciales de 6 ámbitos técnicos (IVA, impuestos especiales, aduanas, clasificación arancelaria, derechos de los pasajeros y comercio de derechos de emisión)\nocite{TJUE2024Prejudiciales}.\footnote{El Tribunal de la Función Pública, creado en 2004, se suprimió en 2016 y sus asuntos pasaron al Tribunal General.}
  \end{itemize}
  Los miembros los nombran los gobiernos de común acuerdo, tras el dictamen de un comité de expertos (art. 255 TFUE), por 6 años renovables.
\item \emph{Funciones} (las principales vías):
  \begin{itemize}
  \item \concepto{cuestión prejudicial} (art. 267 TFUE): el juez nacional pregunta cómo se interpreta una norma de la Unión o si es válida; la respuesta le vincula. Es la vía de la mayoría de las grandes sentencias y la que garantiza una aplicación uniforme;
  \item \emph{recurso por incumplimiento} (arts. 258-260): contra un Estado, con posibles multas;
  \item \emph{recursos de anulación} (art. 263) y \emph{por omisión} (art. 265): control de la legalidad de los actos de las instituciones, como un tribunal constitucional;
  \item \emph{responsabilidad extracontractual} de la Unión (art. 340) y \emph{dictámenes} sobre la compatibilidad de los acuerdos internacionales con los Tratados (art. 218.11).
  \end{itemize}
\end{itemize}

% --------------------------------------------------------------------------
\subsection{El Tribunal de Cuentas}\label{sec:tcu}

\begin{itemize}
\item \emph{Qué es} (art. 285 TFUE): el auditor externo de las finanzas de la Unión. Lo creó el Tratado de Bruselas de 1975, funciona desde 1977 en Luxemburgo y es institución desde Maastricht. No es un tribunal: no juzga ni sanciona.
\item \emph{Composición}: un miembro por Estado, que nombra el Consejo previa consulta al Parlamento, por 6 años renovables; eligen a su presidente para 3 años.
\item \emph{Funciones}:
  \begin{itemize}
  \item examina la \emph{legalidad y regularidad} de los ingresos y gastos y la \emph{buena gestión financiera};
  \item publica el \emph{informe anual} sobre el ejercicio anterior, con la \emph{declaración de fiabilidad} de las cuentas, que el Parlamento usa para aprobar el descargo de la Comisión, e informes especiales;
  \item si detecta fraude, lo comunica a la Oficina Europea de Lucha contra el Fraude (OLAF) o a la Fiscalía Europea \vertema{3.B.37}.
  \end{itemize}
\end{itemize}

% --------------------------------------------------------------------------
\subsection{El Banco Central Europeo}\label{sec:bce}

\begin{itemize}
\item \emph{Qué es} (arts. 282-284 TFUE): el banco central de la zona del euro, con sede en Fráncfort. Lo creó Maastricht, funciona desde junio de 1998 y es institución desde Lisboa. Con los bancos centrales nacionales de los 27 forma el \emph{Sistema Europeo de Bancos Centrales}, y con los de la zona del euro, el \emph{Eurosistema}. Es \emph{independiente}: no puede pedir ni recibir instrucciones.
\item \emph{Objetivo}: la \emph{estabilidad de precios} (concretada en una inflación del 2~\% a medio plazo, simétrica, según su estrategia de 2021, confirmada en 2025) y, sin perjuicio de ella, apoyar las políticas económicas generales de la Unión.
\item \emph{Órganos}:
  \begin{itemize}
  \item el \emph{Comité Ejecutivo}: presidenta (\autor{Christine Lagarde}, que dejará el cargo el 7 de febrero de 2027), vicepresidente (\autor{Boris Vujčić} desde el 1 de junio de 2026, tras \autor{Luis de Guindos})\nocite{HNB2026Vujcic} y 4 miembros, con mandatos de 8 años no renovables que nombra el Consejo Europeo;
  \item el \emph{Consejo de Gobierno}: el Comité Ejecutivo y los 21 gobernadores de la zona del euro (Bulgaria entró el 1 de enero de 2026), que votan por rotación; decide la política monetaria;
  \item el \emph{Consejo General}: presidenta, vicepresidente y los 27 gobernadores, para la relación con los países de fuera del euro.
  \end{itemize}
\item Desde 2014 supervisa además los grandes bancos de la zona del euro (Mecanismo Único de Supervisión). El detalle de sus objetivos, órganos e instrumentos es materia de otros temas \vertema{3.B.43} \vertema{3.B.44}.
\end{itemize}

% --------------------------------------------------------------------------
\subsection{Otros órganos y organismos}\label{sec:otros}

\begin{itemize}
\item \emph{Órganos consultivos} (arts. 300-307 TFUE), que emiten dictámenes obligatorios en los casos que fijan los Tratados y otros por iniciativa propia:
  \begin{itemize}
  \item el \concepto{Comité Económico y Social Europeo} (CESE, 1958): 329 representantes de empresarios, trabajadores y otras organizaciones de la sociedad civil (21 de España);
  \item el \concepto{Comité de las Regiones} (CdR, Maastricht): 329 representantes de entes regionales y locales (21 de España). Puede recurrir al Tribunal por vulneración de la subsidiariedad.
  \end{itemize}
\item El \concepto{Banco Europeo de Inversiones} (BEI; Tratado de Roma, 1958; Luxemburgo): el banco de la Unión, cuyos accionistas son los 27 Estados; financia proyectos de cohesión, clima, innovación y, cada vez más, seguridad y defensa. Lo preside \autor{Nadia Calviño} desde 2024 \vertema{3.B.30}.
\item El \concepto{Defensor del Pueblo Europeo}: lo elige el Parlamento para la legislatura e investiga la mala administración de las instituciones.
\item La \concepto{Fiscalía Europea} (cooperación reforzada; Luxemburgo; actúa desde el 1 de junio de 2021): investiga y lleva a juicio los delitos contra el presupuesto de la Unión (fraude, corrupción, fraude del IVA transfronterizo). Hungría se sumó en julio de 2026 y ya son 25 los Estados participantes (fuera, Dinamarca e Irlanda)\nocite{Eucrim2026Hungria}; desde el 1 de noviembre de 2026 la dirige \autor{Andrés Ritter}, tras \autor{Laura Kövesi}\nocite{EPPO2026Ritter}.
\item El \concepto{Servicio Europeo de Acción Exterior} (2011), el servicio diplomático de la Unión bajo la Alta Representante \vertema{3.B.42}.
\item \emph{Las agencias}: más de 30 agencias descentralizadas con personalidad jurídica propia, que ejecutan tareas técnicas o de supervisión (la del Medicamento, en Ámsterdam; la Autoridad Bancaria Europea, en París; la de Lucha contra el Blanqueo de Capitales, en Fráncfort desde 2025; Frontex, en Varsovia) y agencias ejecutivas de la Comisión que gestionan programas.
\item \emph{En la zona del euro}:
  \begin{itemize}
  \item el \concepto{Eurogrupo}: reunión informal de los ministros de Finanzas de la zona del euro (Protocolo 14), con un presidente elegido para 2 años y medio; desde el 12 de diciembre de 2025 es el griego \autor{Kyriakos Pierrakakis}\nocite{Eurogrupo2025Presidente};
  \item el \concepto{Mecanismo Europeo de Estabilidad} (MEDE, 2012): no es un órgano de la Unión, sino una organización intergubernamental de los Estados del euro, con sede en Luxemburgo, que presta a los que pierden el acceso a los mercados con condicionalidad, con una capacidad de préstamo de 500.000 millones de euros. La reforma de su tratado (2021) no ha entrado en vigor porque Italia no la ha ratificado \vertema{3.B.45}.
  \end{itemize}
\end{itemize}

\begin{anotaciones}
\textbf{3 «consejos» que no hay que confundir.}
\begin{itemize}
\item El \emph{Consejo Europeo}: los jefes de Estado o de Gobierno de la Unión; da el impulso político (\ref{sec:consejo-europeo}).
\item El \emph{Consejo de la Unión Europea} (o Consejo): los ministros de los Estados; legisla con el Parlamento (\ref{sec:consejo}).
\item El \emph{Consejo de Europa}: una organización internacional distinta, de 46 Estados, con sede en Estrasburgo, que protege los derechos humanos con el Convenio Europeo y su Tribunal. La Unión está obligada a adherirse al Convenio (art. 6.2 TUE), pero el Tribunal de Justicia rechazó el primer acuerdo de adhesión en su dictamen 2/13 (2014) y la negociación se reabrió después.
\end{itemize}
\end{anotaciones}

\begin{itemize}
\item \emph{Recapitulación}: las instituciones reparten el poder entre 3 legitimidades, la de los ciudadanos (Parlamento), la de los Estados (Consejo y Consejo Europeo) y la del interés común (Comisión, Tribunal de Justicia y BCE). Las reglas de voto del Consejo deciden cuánto pesa cada Estado, y la tendencia de cada reforma ha sido más mayoría y más Parlamento.
\end{itemize}

% ==========================================================================
\section*{Conclusión}
\minutos{1}

\begin{itemize}
\item \textbf{\emph{Recapitulación}}:
  \begin{itemize}
  \item En conclusión, la Unión se construyó como anunció \autor{Schuman}: por realizaciones concretas. De París a Lisboa, cada tratado convirtió una ambición en una regla, y la Unión pasó de 6 a 27 Estados y del carbón y el acero a una moneda común.
  \item Su Derecho es lo que la distingue de cualquier otra organización: prima sobre el nacional, se aplica directamente y lo garantiza un tribunal propio. La subsidiariedad reparte las competencias con una lógica económica: centralizar donde las externalidades superan a la diversidad de preferencias.
  \item Sus instituciones equilibran a los ciudadanos, a los Estados y al interés común, y sus reglas de voto reparten un poder que no es proporcional al peso.
  \end{itemize}
\item \textbf{\emph{Relevancia}}:
  \begin{itemize}
  \item Para España, la Unión es el marco de casi toda su política económica: comercial, monetaria, de competencia, agrícola y presupuestaria, y una fuente de fondos y de reformas \vertema{4.A.30}.
  \end{itemize}
\item \textbf{\emph{Extensiones}}:
  \begin{itemize}
  \item La integración económica en teoría \vertema{3.B.10} \vertema{3.B.16}; el presupuesto y el marco financiero \vertema{3.B.37}; el mercado interior \vertema{3.B.39}; la política comercial \vertema{3.B.41}; el euro y su gobernanza \vertema{3.B.43} \vertema{3.B.45}; y el federalismo fiscal \vertema{4.B.14}.
  \end{itemize}
\item \textbf{\emph{Opinión}}:
  \begin{itemize}
  \item La ampliación hacia el Este y los Balcanes es la gran decisión de esta década. Hacerla con las reglas actuales arriesga la parálisis; esperar a una reforma de los Tratados arriesga no hacerla nunca. Lo realista es lo que se ha empezado a hacer en 2026: usar las pasarelas y la cooperación reforzada, ligar la entrada a salvaguardias de los valores y dejar para más adelante la reforma completa.
  \end{itemize}
\item \textbf{\emph{Idea final}}:
  \begin{itemize}
  \item Volvamos al refrán: del dicho al hecho hay mucho trecho. La historia de la Unión es la de ese trecho recorrido a base de Derecho: cada tratado convirtió un dicho en un hecho jurídico. Sus reglas le han servido para crecer de 6 a 27; para llegar a más de 30, el trecho que queda es, otra vez, el de convertir en reglas lo que ya se ha prometido.
  \end{itemize}
\end{itemize}

\clearpage
\appendix

% ==========================================================================
\section{Anexo: los Estados miembros}\label{anexo:estados}

\begin{table}[H]
\centering
\caption{Los 27 Estados miembros: adhesión, euro, Schengen, escaños y población}\label{tab:estados}
\small
\begin{tabular}{lcccrr}
\hline
\textbf{Estado} & \textbf{Adhesión} & \textbf{Euro} & \textbf{Schengen} & \textbf{Escaños} & \textbf{Población (\%)}\\
\hline
Alemania & 1952 & 1999 & Sí & 96 & 18,47\\
Francia & 1952 & 1999 & Sí & 81 & 15,18\\
Italia & 1952 & 1999 & Sí & 76 & 13,19\\
España & 1986 & 1999 & Sí & 61 & 10,86\\
Polonia & 2004 & --- & Sí & 53 & 8,29\\
Rumanía & 2007 & --- & Sí (2025) & 33 & 4,21\\
Países Bajos & 1952 & 1999 & Sí & 31 & 4,02\\
Bélgica & 1952 & 1999 & Sí & 22 & 2,63\\
Chequia & 2004 & --- & Sí & 21 & 2,42\\
Portugal & 1986 & 1999 & Sí & 21 & 2,38\\
Suecia & 1995 & --- & Sí & 21 & 2,34\\
Grecia & 1981 & 2001 & Sí & 21 & 2,30\\
Hungría & 2004 & --- & Sí & 21 & 2,11\\
Austria & 1995 & 1999 & Sí & 20 & 2,03\\
Bulgaria & 2007 & 2026 & Sí (2025) & 17 & 1,42\\
Dinamarca & 1973 & --- & Sí & 15 & 1,32\\
Finlandia & 1995 & 1999 & Sí & 15 & 1,24\\
Eslovaquia & 2004 & 2009 & Sí & 15 & 1,21\\
Irlanda & 1973 & 1999 & No & 14 & 1,20\\
Croacia & 2013 & 2023 & Sí (2023) & 12 & 0,86\\
Lituania & 2004 & 2015 & Sí & 11 & 0,64\\
Eslovenia & 2004 & 2007 & Sí & 9 & 0,47\\
Letonia & 2004 & 2014 & Sí & 9 & 0,41\\
Estonia & 2004 & 2011 & Sí & 7 & 0,30\\
Chipre & 2004 & 2008 & No & 6 & 0,22\\
Luxemburgo & 1952 & 1999 & Sí & 6 & 0,15\\
Malta & 2004 & 2008 & Sí & 6 & 0,13\\
\hline
\end{tabular}
\fuente{elaboración propia a partir de la Unión Europea\nocite{UE2026Instituciones}, la Decisión 2023/2061 del Consejo Europeo\nocite{ConsejoEuropeo2023Composicion} y la Decisión 2025/2427 del Consejo\nocite{Consejo2025Poblacion}. Adhesión de los fundadores: CECA (1952). Bulgaria y Rumanía, en Schengen desde el 1 de enero de 2025 (por aire y mar desde el 31 de marzo de 2024).}
\end{table}

% ==========================================================================
\section{Anexo: el cálculo de los índices de poder}\label{anexo:poder}

\begin{itemize}
\item \emph{El problema}: con 27 Estados hay $2^{26}\approx6{,}7\cdot10^{7}$ coaliciones de los demás para cada Estado; contarlas una a una es posible, pero lento. Se puede hacer exacto con \emph{programación dinámica}, porque la regla solo depende de 2 números de la coalición: cuántos Estados tiene ($k$) y cuánta población ($p$).
\item \emph{Paso 1}: para cada Estado $i$, se cuenta cuántas coaliciones de los demás hay con cada par $(k,p)$, $C_i(k,p)$, añadiendo los Estados de uno en uno: al añadir el Estado $j$, de población $w_j$, $C(k,p)\leftarrow C(k,p)+C(k-1,p-w_j)$.
\item \emph{Paso 2}: $i$ es decisivo en las coaliciones $(k,p)$ que pierden sin él y ganan con él; con la regla $v=1\iff(k\geq15\ \text{y}\ p\geq0{,}65\,P)\ \text{o}\ k\geq24$, son las que cumplen
  \begin{itemize}
  \item $k+1\geq15$ y $p+w_i\geq0{,}65\,P$, pero no $k\geq15$ y $p\geq0{,}65\,P$; o bien
  \item $k=23$ (con él llegan a 24 Estados).
  \end{itemize}
\item \emph{Paso 3}: el índice de \autor{Banzhaf} suma las coaliciones decisivas y las divide por $2^{26}$; el de \autor{Shapley}-\autor{Shubik} las pondera por $k!\,(26-k)!/27!$, la probabilidad de que en un orden aleatorio lleguen antes que $i$ justo esos $k$ Estados.
\item \emph{Resultado} (población en miles, para que el cálculo sea manejable; tabla~\ref{tab:poder}). La probabilidad de que una coalición al azar gane es del 13,2~\%.\findemo
\end{itemize}

\begin{table}[H]
\centering
\caption{Población y poder de voto a priori en el Consejo, 2026 (\%)}\label{tab:poder}
\small
\begin{tabular}{lrrr@{\hspace{1.2cm}}lrrr}
\hline
\textbf{Estado} & \textbf{Pob.} & \textbf{SS} & \textbf{Bz} & \textbf{Estado} & \textbf{Pob.} & \textbf{SS} & \textbf{Bz}\\
\hline
Alemania & 18,47 & 17,54 & 12,03 & Austria & 2,03 & 2,23 & 2,75\\
Francia & 15,18 & 13,48 & 10,12 & Bulgaria & 1,42 & 1,79 & 2,43\\
Italia & 13,19 & 11,53 & 8,87 & Dinamarca & 1,32 & 1,73 & 2,38\\
España & 10,86 & 9,39 & 7,81 & Finlandia & 1,24 & 1,67 & 2,34\\
Polonia & 8,29 & 6,69 & 6,32 & Eslovaquia & 1,21 & 1,64 & 2,32\\
Rumanía & 4,21 & 3,79 & 3,90 & Irlanda & 1,20 & 1,64 & 2,31\\
Países Bajos & 4,02 & 3,66 & 3,80 & Croacia & 0,86 & 1,38 & 2,13\\
Bélgica & 2,63 & 2,65 & 3,07 & Lituania & 0,64 & 1,23 & 2,01\\
Chequia & 2,42 & 2,50 & 2,96 & Eslovenia & 0,47 & 1,11 & 1,92\\
Portugal & 2,38 & 2,47 & 2,94 & Letonia & 0,41 & 1,07 & 1,89\\
Suecia & 2,34 & 2,45 & 2,92 & Estonia & 0,30 & 0,99 & 1,83\\
Grecia & 2,30 & 2,42 & 2,90 & Chipre & 0,22 & 0,93 & 1,78\\
Hungría & 2,11 & 2,28 & 2,80 & Luxemburgo & 0,15 & 0,88 & 1,75\\
& & & & Malta & 0,13 & 0,87 & 1,73\\
\hline
\end{tabular}
\fuente{elaboración propia con las cifras de población de 2026\nocite{Consejo2025Poblacion}. SS: \autor{Shapley}-\autor{Shubik}; Bz: \autor{Banzhaf} normalizado.}
\end{table}

% ==========================================================================
\section{Anexo: preguntas de test}\label{anexo:test}

\begin{anotaciones}
Preguntas tipo test del primer ejercicio: \href{https://www.victorgutierrezmarcos.es/oposicion/temario/primer-ejercicio/test/simulador.html?tema=3.B.36}{simulador de test}.
\end{anotaciones}

\bibliografia
\end{document}
